\providecommand{\toplevelprefix}{../..}  %
\documentclass[../../book-main_zh.tex]{subfiles}

\begin{document}

\chapter{关于智能的非正式引言}
\label{ch:intro}

\begin{quote}
    ``{\em 正如熵的不断增加是宇宙的基本定律，
        生命的基本定律则是变得越来越
    高度结构化并与熵作斗争。}''

    $~$\hfill -- 瓦茨拉夫·哈维尔 (V\'{a}clav Havel)
\end{quote}
\vspace{5mm}

\section{智能、控制论与人工智能}
我们居住的世界既不是完全随机的，也不是完全
不可预测的。\footnote{如果世界是完全随机的，一个
智能体就没有必要去学习或记忆任何东西。}
相反，它遵循着某些秩序、模式和规律，这使得它
在很大程度上是可预测的。\footnote{有些是确定性的，有些是
概率性的。} 生命的出现和延续正是依赖于
这种可预测性。只有通过学习和记忆环境中
可预测的事物，生命才能生存和繁衍，因为
合理的决策和行动取决于可靠的预测。由于
世界提供了似乎无限的可预测现象，
智能生物——动物和人类——进化出了越来越
敏锐的感官：视觉、听觉、触觉、味觉和嗅觉。这些感官
收集高通量的感觉数据以感知环境的
规律性。因此，所有智能生物的一个基本任务是
\begin{center}
    \textit{从海量的感知数据中学习并记忆可预测的
    信息}。
\end{center}
在理解这是如何实现之前，我们必须解决
三个问题：
\begin{itemize}
    \item 如何对可预测的信息进行数学建模和
        表征？
    \item 如何通过计算从数据中有效且高效地
        学习这些信息？
    \item 这些信息应该如何以最佳方式被组织和结构化为记忆或（经验）知识，存储在我们的大脑或机器中，以支持未来的预测和推理？
\end{itemize}
本书旨在为这些问题提供一些答案。这些
答案将帮助我们更好地理解智能，特别是
实现智能的计算原理和机制。有证据
表明，所有形式的智能——从早期原始生命中可见的低级智能
到最高形式的智能，即
现代科学的实践——都共享着一套共同的原理和
机制。我们在下文进行详细阐述。

\paragraph{智能的出现与进化。}

大约四十亿年前地球上生命出现的一个必要条件是
环境在很大程度上是可预测的。
生命已经发展出一些机制，使其能够学习环境中
可预测的事物，对这些信息进行编码，并将其用于
生存。一般来说，我们将这种学习世界知识的能力称为
\textit{智能}。在很大程度上，
生命的进化就是智能机制在发挥作用
\cite{Bennett-2023}。在生命的早期阶段，智能主要
通过两种类型的学习机制发展而来：
\textit{系统发育（phylogenetic）}和\textit{个体发育（ontogenetic）}
\cite{Wiener-Cybernetics-1961}。

\textit{系统发育智能}是指通过
物种的进化进行学习。物种的遗传和生存主要基于编码在
其父母的DNA或基因中的记忆和
知识。在很大
程度上，我们可以将DNA称为大自然的预训练大模型，因为
它们扮演着非常相似的角色。
系统发育智能的主要特征是个体的学习能力有限。
学习是通过基于基因随机突变的“试错”
机制进行的，而物种的进化则基于
自然选择——适者生存——如
\Cref{fig:phylogenetic}所示。这可以被视为大自然对
现在所谓的“强化学习”的实现，
\cite{Sutton-Barto} 或潜在的“神经网络架构搜索”
\cite{NAS-1}。然而，这种“试错”过程可能
极其缓慢、昂贵且不可预测。从大约44至38亿年前
第一批生命形式出现以来
\citep{bell2015potentially}，生命一直依赖于这种形式的
进化。\footnote{敏锐的读者可能会注意到，早期生命的进化方式与
当今大型语言模型的进化方式之间存在着惊人的相似性。}

\begin{figure}
    \centering
    \includegraphics[width=0.5\linewidth]{\toplevelprefix/chapters/chapter1/figs/DNAs.png}
    \includegraphics[width=0.40\linewidth]{\toplevelprefix/chapters/chapter1/figs/Evolution.jpg}
    \caption{系统发育智能的进化：外部世界的记忆或知识
        通过DNA进行编码和传递（左），并
        从DNA解码为RNA再到蛋白质。在生命进化的早期阶段
        （右），智能在物种层面上通过（随机的）基因突变和
        自然选择——“适者生存”——来发展记忆或知识，
    这可以被视为强化学习的一种原始形式。}
    \label{fig:phylogenetic}
\end{figure}

\begin{figure}
    \centering
    \includegraphics[height=0.19\linewidth]{\toplevelprefix/chapters/chapter1/figs/Luca.jpeg}
    \includegraphics[height=0.19\linewidth]{\toplevelprefix/chapters/chapter1/figs/Worm.jpeg}
    \includegraphics[height=0.19\linewidth]{\toplevelprefix/chapters/chapter1/figs/Cambrian.jpg}
    \caption{生命的进化，从今天所有生命的祖先
        （LUCA——最后广泛共同祖先），一种生活在35至43亿年前的类似单细胞的
        生物
        \cite{moody2024nature}，到大约5.5亿年前在蠕虫类物种中
        出现的第一个神经系统（中）
        \cite{wang2025preservation}，再到大约5.3亿年前
    寒武纪时期的生命形式大爆发（右）。}
    \label{fig:evolution}
\end{figure}

\textit{个体发育智能}是指允许个体在其特定环境中
通过自身的感官、记忆和
预测进行学习，并改善和
适应其行为的学习机制。在大约5.5至6亿年前（在
蠕虫类生物中）神经系统出现之后，个体发育学习成为可能
\cite{wang2025preservation}，如
\Cref{fig:evolution}中间所示。借助感觉和神经系统，一个
个体除了从DNA或基因中继承的内容外，还可以不断形成和改善其自身关于
世界的记忆。
这种能力显著增强了个体的生存能力，并
促成了大约5.3亿年前寒武纪生命形式的大爆发
\cite{Parker-2004}。与系统发育学习相比，
个体发育学习更高效、更可预测，并且可以
在个体的资源限制内实现。

这两种类型的学习都依赖于来自外部
环境的反馈——惩罚（死亡）或奖励（食物）——这些反馈作用于
物种或个体的行动上。\footnote{物种的基因突变
或个体做出的行动。} 这一见解启发了
诺伯特·维纳（Norbert Wiener）在其控制论（Cybernetics）项目中得出结论
\cite{Wiener-Cybernetics-1948}，所有智能生物，无论是
物种还是个体，都依赖于闭环反馈机制来
学习和改善它们关于世界的记忆和知识。此外，从
植物到鱼类、鸟类和哺乳动物，更高级的物种
越来越依赖于个体发育学习：它们在出生后会与父母待在一起并向父母学习
更长的时间，因为
同一物种的个体必须在非常多样化的环境中生存。

\paragraph{人类智能的进化。}
自从大约250万年前\textit{智人（Homo sapiens）}出现以来
\cite{Harari-2015}，一种新的、更高形式的智能已经
出现，它以更高效、更经济的方式进化。人类
社会发展出了语言——首先是口语，后来是书面语——如
\Cref{fig:human-intelligence}所示。语言使个体能够
交流和分享有用的信息，从而使人类
社区能够像一个单一的智能有机体一样运作，其学习
速度比任何个体都快，保留的知识也比任何个体都多。因此，书面文本
扮演着类似于DNA和基因的角色，使社会能够
跨代积累和传递知识。我们可以将
这种类型的智能称为\textit{社会智能}，
以区别于物种的系统发育智能和
个体的个体发育智能。这种知识的
积累构成了（古代）文明的基础。

\begin{figure}
    \centering
    \includegraphics[height=0.25\linewidth]{\toplevelprefix/chapters/chapter1/figs/Spoken-language.jpg}
    \hspace{5mm}
    \includegraphics[height=0.25\linewidth]{\toplevelprefix/chapters/chapter1/figs/Cuneiform.png}
    \hspace{5mm}
    \includegraphics[height=0.25\linewidth]{\toplevelprefix/chapters/chapter1/figs/adopt-euclid1685-2.jpg}
    \caption{口头交流和口语（7万至3万年前）、
        书面语言（约公元前3000年）
        \cite{schmandt2014evolution}以及抽象数学
        （约公元前500至300年）\cite{heath1956thirteen}的发展标志着
    人类智能进化中的三个关键里程碑。}
    \label{fig:human-intelligence}
\end{figure}

大约两三千年前，人类智能实现了
另一次重大飞跃，使哲学家和数学家能够
发展出远远超出组织经验
观察的知识。抽象概念和符号（如
数字、时间、空间、逻辑和几何）的发展，催生了
一种全新的、严谨的\textit{数学}语言。此
外，提出假设并通过逻辑推演或实验验证其正确性的
能力的发展，
为现代\textit{科学}奠定了基础。人类首次
能够主动、系统地发现和发展新
知识。我们将这种高级形式的智能称为
“科学智能”，因为它是演绎和
科学发现所必需的。

因此，从我们能向大自然学习到的东西来看，每当我们使用
“智能”一词时，我们必须明确我们指的是哪个层次或形式：
\begin{equation}
    \text{\textbf{phylogenetic}} \;
    \Longrightarrow \; \text{\textbf{ontogenetic}} \; \Longrightarrow \;
    \text{\textbf{societal}}
    \; \Longrightarrow \;
    \text{\textbf{scientific intelligence}}.
\end{equation}
清晰的刻画和区分是必要的，因为我们希望
将智能作为一门科学和数学学科来研究。
尽管所有形式可能都共享着学习和
改善关于世界的记忆和知识的共同目标，但计算机制和物理实现可能会有所不同。我们相信读者在
学习本书后，将能更好地理解和体会它们的共性与差异。
因此，我们将关于这些不同形式或层次的智能之间关系的
更深层次的科学或哲学讨论
留到本书的最后\Cref{ch:future}。

\paragraph{机器智能的起源——控制论。}
在20世纪40年代，受战争需求的推动，受到
自然智能启发的科学家们试图用机器来模拟动物智能，
从而催生了由诺伯特·维纳倡导的“控制论”运动
\cite{Kline-2011}。维纳在成为数学家和控制理论家之前，
曾在哈佛大学本科学习动物学。
他毕生致力于理解和构建
能够重现类似动物智能的自主系统。
今天，控制论项目通常被狭隘地解释为主要
关于反馈控制系统，这是维纳做出
最重大技术贡献的领域。然而，该项目要
广泛和深刻得多：它旨在将智能作为一个
整体来理解——至少在动物层面上——并影响了
整整一代著名科学家的工作，包括沃伦·麦卡洛克（Warren McCulloch）、
沃尔特·皮茨（Walter Pitts）、克劳德·香农（Claude Shannon）、约翰·冯·诺依曼（John von
Neumann）和艾伦·图灵（Alan Turing）。

\begin{figure}
    \centering
    \includegraphics[height=0.4\linewidth]{\toplevelprefix/chapters/chapter1/figs/Cybernetics1.jpg}
    \hspace{10mm}
    \includegraphics[height=0.4\linewidth]{\toplevelprefix/chapters/chapter1/figs/Cybernetics2.jpg}
    \caption{诺伯特·维纳的著作《控制论》（1948年）
        \cite{Wiener-Cybernetics-1948}（左）及其第二版
    （1961年）\cite{Wiener-Cybernetics-1961}（右）。}
    \label{fig:cybernetcis}
\end{figure}

维纳可以说是第一个\textit{作为一个系统}来研究智能的人，
而不是专注于智能的孤立组件或方面。
他的全面观点出现在1948年著名的
著作《控制论：或关于在动物和机器中控制和通信的科学》
\cite{Wiener-Cybernetics-1948}中。在该书及其
1961年出版的第二版\cite{Wiener-Cybernetics-1961}（见
\Cref{fig:cybernetcis}）中，他试图确定智能系统的几个必要
特征和机制，包括（但
不限于）：
\begin{itemize}
    \item 如何（在大脑中）\textit{测量和存储}信息
        以及如何与他人交流。\footnote{维纳
            是第一个指出“信息”既不是
            物质也不是能量，而是一个值得研究的独立量的人。
        } 这一见解促使克劳德·香农在1948年提出了
        \textit{信息论} \cite{Shannon-1948}。
    \item 如何基于现有信息\textit{纠正}预测和估计中的\textit{错误}
        。维纳本人在20世纪40年代帮助
        形式化了（闭环）反馈控制理论。
    \item 在与不合作甚至对抗性的环境互动时，如何学习\textit{做出更好的决策}
        。
        约翰·冯·诺依曼在1944年将其形式化为\textit{博弈论}
        \cite{von-Neumann-1944}。
\end{itemize}
1943年，受维纳控制论项目的启发，认知
科学家沃伦·麦卡洛克和逻辑学家沃尔特·皮茨共同
形式化了第一个神经元的计算模型
\cite{McCulloch-Pitts}，称为\textit{人工神经元}，
稍后在\Cref{fig:neuron}中进行了说明。基于该模型，弗兰克·罗森布拉特（Frank Rosenblatt）
在20世纪50年代构建了Mark I感知机——一台
包含数百个此类人工神经元的物理机器
\cite{Rosenblatt:1957}。感知机是第一个物理
实现的人工神经网络；见\Cref{fig:perceptron}。
值得注意的是，约翰·冯·诺依曼在1945年提出的通用计算机架构，
也是为了促进构建
能够物理实现控制论项目所建议机制的\textit{计算机器}
这一目标而设计的 \cite{von-Neumann-1958}。

敏锐的读者会注意到，20世纪40年代确实是神奇的：许多
基本思想被发明出来，有影响力的理论被形式化，
包括神经元的数学模型、人工神经
网络、信息论、控制理论、博弈论和
计算机器。\Cref{fig:god-fathers}描绘了其中的一些
先驱。这些贡献中的每一项都已发展成为
科学或工程领域的基础，并继续对
我们的生活产生巨大影响。所有这些都是受到开发
模拟自然界智能的机器这一目标的启发。历史
记录表明，维纳的控制论运动几乎影响了所有
这些先驱和工作。在很大程度上，维纳的项目可以
被视为当今“具身
智能”项目的真正前身；事实上，维纳本人以极其清晰和具体的方式
阐述了这样一个项目的愿景
\cite{Wiener-Cybernetics-1961}。
\begin{figure}
    \centering
    \includegraphics[height=0.3\linewidth]{\toplevelprefix/chapters/chapter1/figs/Wiener.png}
    \includegraphics[height=0.3\linewidth]{\toplevelprefix/chapters/chapter1/figs/Shannon.jpg}
    \includegraphics[height=0.3\linewidth]{\toplevelprefix/chapters/chapter1/figs/neumann.jpg}
    \includegraphics[height=0.3\linewidth]{\toplevelprefix/chapters/chapter1/figs/Turing.jpeg}
    \caption{智能理论和计算基础的先驱：
        诺伯特·维纳（控制论和控制
        理论）、克劳德·香农（信息论）、约翰·冯·诺依曼
    （博弈论）和艾伦·图灵（计算理论）。}
    \label{fig:god-fathers}
\end{figure}

尽管维纳确定了（具身）智能的许多关键特征和机制，
但他没有提供将它们整合到
一个完整的自主智能系统中的明确方法。从今天的
角度来看，他的一些观点是不完整或不准确的。在
特别是，在《控制论》第二版的最后一章中
\cite{Wiener-Cybernetics-1961}，他强调了
如果机器要模拟自然界中典型的学习机制，就需要\textit{处理非线性}，
然而他并没有为这个难题提供具体、
有效的解决方案。平心而论，即使是
线性系统理论在当时也处于起步阶段，而
非线性系统似乎更难以企及。

尽管如此，我们还是不禁对维纳关于
非线性重要性的先见之明感到惊叹。正如本书将展示的那样，答案
直到最近才出现：非线性可以通过
深度网络和表征实现的渐进式线性化和变换来有效处理
（见\Cref{ch:deep-networks}）。
此外，我们将在本书中展示如何将上述所有机制
自然地整合到一个完整的系统中，该系统
表现出自主智能系统的某些特征
（见\Cref{ch:consistent-representations}）。

\paragraph{人工智能的起源。}
维纳《控制论》的副标题——\textit{在动物和机器中控制和
通信的科学}——揭示了20世纪40年代的
研究主要旨在模拟动物级别的智能。正如
前面提到的，那个时代的议程由维纳的
控制论运动主导。

艾伦·图灵是最早认识到这一局限性的人之一。在他
著名的1950年论文《\textit{计算机器与
智能}》\cite{Turing-1950}中，图灵正式提出了
机器是否能够模仿人类级别智能的问题，
以至于机器智能与人类智能
无法区分——现在被称为\textit{图灵测试}。

大约在1955年，一群雄心勃勃的年轻科学家试图打破
当时占主导地位的控制论项目，并建立他们
自己的遗产。他们接受了图灵关于模仿人类
智能的挑战，并提议于1956年夏天在达特茅斯学院举办
一次研讨会。他们的提案中写道\cite{McCarthy-1955}：
\begin{quote}
    ``\textit{该研究基于这样一个猜想，即学习的每一个方面
        或智能的任何其他特征，
        原则上都可以被精确地描述，以至于可以
        制造一台机器来模拟它。我们将尝试
        寻找如何让机器使用语言、形成抽象和
        概念、解决目前只有人类才能解决的各种问题，以及
    自我改进}。''
\end{quote}
他们旨在形式化和研究区分人类与动物的
更高层次的智能。他们的议程包括抽象、
符号方法、自然语言和演绎推理（因果
推理、逻辑演绎等），即科学智能。
研讨会的组织者约翰·麦卡锡（John McCarthy）——当时是达特茅斯学院的一名年轻数学助理
教授——创造了现在著名的
术语“人工智能”（AI）来描述表现出
科学智能的机器，以便正式描述
研讨会的问题和目标。

\paragraph{“人工智能”还是“控制论”的复兴？}
在过去的十年里，机器智能经历了爆炸式的
发展，这在很大程度上是由深度人工神经网络驱动的，
并由杰弗里·辛顿（Geoffrey Hinton）及其学生在2012年的工作所引发
\cite{krizhevsky2012imagenet}。这一时期被誉为
“人工智能的复兴”。然而，就实际解决的任务
（识别、生成、预测）和开发的技术
（强化学习、模仿学习、编码、解码、
去噪和压缩）而言，我们在很大程度上是在模拟
早期生命和动物智能所共有的机制。即使在这
方面，正如我们将在本书中阐明的那样，当前的“AI”模型和
系统也没有完全或正确地实现20世纪40年代控制论运动所知的
动物级别智能的所有必要
机制。

严格来说，过去十年机器智能的最新进展
与1956年达特茅斯人工智能项目并不十分吻合。
相反，所取得的成就更接近于
维纳经典的20世纪40年代控制论项目的目标。将当前时代称为
“控制论的复兴”可能更为
合适。最近兴起的用于自主机器人和智能体的所谓“具身AI”
甚至更紧密地契合了
控制论项目的目标。只有当我们从科学和
数学的角度完全理解我们真正取得了什么成就之后，我们才能
确定哪些问题仍然悬而未决，以及哪些方向通向智能的
真正本质
。\footnote{因此，我们将在本书的最后
一章\ref{ch:future}中恢复我们的讨论。} 这是本书的
主要目的之一。

\section{学习什么？}
\label{sec:what-to-learn}

\subsection{可预测性}
\label{sec:predictability}
携带有用信息的数据以多种形式表现出来。在它们
最自然的形式中，它们可以被建模为可预测和
可计算的序列。可预测和
可计算序列的概念是计算理论的核心，并
在很大程度上导致了计算机的发明 \cite{Turing-1936}。
可预测序列在（归纳）推理中的作用由
雷·所罗门诺夫（Ray Solomonoff）、安德雷·柯尔莫哥洛夫（Andrey Kolmogorov）等人在20世纪60年代进行了研究
\cite{Kolmogorov1998OnTO}，作为克劳德·香农
经典信息论的推广 \cite{Shannon-1948}。为了理解
可预测序列的概念，我们从一些具体的例子开始。

\paragraph{标量情况。} 最简单的可预测离散序列
可以说是自然数序列：
\begin{equation}
    S = 1, 2, 3, 4, 5, 6, \ldots, n, n+1, \ldots
\end{equation}
其中下一个数字 \(x_{n+1}\) 被定义为前一个
数字 \(x_{n}\) 加1：
\begin{equation}
    x_{n+1} = x_{n} + 1.
\end{equation}
人们可以将可预测性的概念推广到任何序列
\(\{x_{n}\}_{n=1}^{\infty}\)（其中 \(x_{n} \in \R\)），如果下一个数字
\(x_{n+1}\) 总是可以从其前驱 \(x_{n}\) 计算得出：
\begin{equation}
    x_{n+1} = f(x_{n}), \quad x_{n} \in \R, \; n = 1, 2, 3, \ldots
\end{equation}
其中 \(f(\cdot)\) 是一个\textit{可计算的}（标量）
函数。\footnote{这里我们强调函数 \(f(\cdot)\)
    本身是可计算的，这意味着它可以作为程序在
计算机上实现。} 艾伦·图灵在1936年的开创性工作 \cite{Turing-1936}
给出了可计算性的严格定义。在实践中，我们通常
进一步假设 \(f\) 是可高效计算的，并且具有良好的
性质，如连续性和可微性。一旦我们理解了
更精细的可计算性概念及其在机器学习和智能中的作用，
这些性质的必要性
在稍后就会变得清晰。

\paragraph{多变量情况。}
下一个值也可能依赖于两个前驱。例如，著名的
\textit{斐波那契数列}
\begin{equation}
    S = 1, 1, 2, 3, 5, 8, 13, 21, 34, 55, \ldots
\end{equation}
满足
\begin{equation}
    x_{n+2} = x_{n+1} + x_{n}, \quad x_{n} \in \R, \; n = 1, 2, 3, \ldots
\end{equation}
更一般地，我们可以写成
\begin{equation}
    x_{n+2} = f(x_{n+1}, x_{n}), \quad x_{n} \in \R, \; n = 1, 2, 3, \ldots
\end{equation}
对于任何接受两个输入的可计算函数 \(f\)。进一步扩展，
下一个值可能依赖于前面的 \(d\) 个值：
\begin{equation}
    x_{n+d} = f(x_{n+d-1}, \ldots, x_{n}), \quad x_{n} \in \R, \; n =
    1, 2, 3, \ldots
    \label{eqn:recursive-d}
\end{equation}
整数 \(d\) 被称为递归的\textit{阶数}。
上述表达式 \eqref{eqn:recursive-d} 被称为
\textit{自回归}，由此产生的序列是
\textit{自回归的}。当 \(f\) 是线性时，我们称其为
\textit{线性自回归}。

\paragraph{向量情况。}
为了简化符号，我们定义一个向量 \(\vx \in \R^{d}\)，它
收集了序列中连续的 \(d\) 个值：
\begin{equation}
    \vx_{n} \doteq [x_{n+d-1}, \ldots, x_{n}]^{\top}, \quad \vx_{n}
    \in \R^{d}, \; n = 1, 2, 3, \ldots
\end{equation}
使用这种符号，递归关系 \eqref{eqn:recursive-d} 变为
\begin{equation}
    \vx_{n+1} = g(\vx_{n}) \in \R^{d}, \quad n = 1, 2, 3, \ldots
    \label{eqn:recursive-v}
\end{equation}
其中 \(g(\cdot)\) 由 \eqref{eqn:recursive-d} 中的函数 \(f\)
唯一确定，并将一个 \(d\) 维向量映射到一个
\(d\) 维向量。在不同的上下文中，这样的向量有时
被称为“状态（state）”或“标记（token）”。请注意，
\eqref{eqn:recursive-d} 定义了一个映射 \(\R^{d} \to \R\)，而
这里我们有 \(g \colon \R^{d} \to \R^{d}\)。

\paragraph{受控预测。}
我们还可以定义一个依赖于另一个
可预测序列作为输入的可预测序列：
\begin{equation}
    \vx_{n+1} = f(\vx_{n}, \vu_{n}) \in \R^{d}, \quad n = 1, 2, 3, \ldots,
    \label{eqn:recursive-control}
\end{equation}
其中 \(\{\vu_{n}\}\)（\(\vu_{n} \in \R^{k}\)）是一个（可计算的）
可预测序列。换句话说，下一个向量 \(\vx_{n+1} \in
\R^{d}\) 既依赖于 \(\vx_{n} \in \R^{d}\) 也依赖于 \(\vu_{n} \in
\R^{k}\)。在控制理论中，序列 \(\{\vu_{n}\}\) 通常
被称为“控制输入”，而 \(\vx_{n}\) 被称为系统 \eqref{eqn:recursive-control} 的“状态”
或“输出”。一个经典的
例子是线性动力系统：
\begin{equation}
    \vx_{n+1} = \vA\vx_{n} + \vB\vu_{n}, \quad \vA \in \R^{d \times
    d}, \vB \in \R^{d \times k},
    \label{eqn:lineary-system}
\end{equation}
这在控制理论中得到了广泛的研究 \cite{Cal:Des}。

通常，控制输入由状态 \(\vx_{n}\) 本身的
可计算函数给出：
\begin{equation}
    \vu_{n} = h(\vx_{n}), \quad n = 1, 2, 3, \ldots
    \label{eqn:control-function}
\end{equation}
因此，序列 \(\{\vx_{n}\}\) 是由两个可计算函数 \(f\) 和 \(h\) 复合而成的：
\begin{equation}
    \vx_{n+1} = f(\vx_{n}, h(\vx_{n})), \quad n = 1, 2, 3, \ldots
    \label{eqn:recursive-closed-loop}
\end{equation}
这样一来，序列 \(\{\vx_{n}\}\) 再次成为一个自回归的可预测序列。当输入 \(\vu_{n}\)
依赖于输出 \(\vx_{n}\) 时，我们称由此产生的序列是由“闭环”系统
\eqref{eqn:recursive-closed-loop} 产生的。由于闭环系统不再
依赖于任何外部输入，我们称这样的系统已经变得\textit{自治}（autonomous）。
它可以被视为自回归的一个特例。例如，如果我们在上述线性系统 \eqref{eqn:lineary-system} 中
选择 \(\vu_{n} = \vF\vx_{n}\)，
闭环系统将变为
\begin{equation}
    \vx_{n+1} = \vA\vx_{n} + \vB\vu_{n} = \vA\vx_{n} + \vB\vF\vx_{n}
    = (\vA + \vB\vF)\vx_{n},
    \label{eqn:lineary-system-closed}
\end{equation}
这是一个线性自回归。

\paragraph{连续过程。}
可预测序列具有自然的连续对应物，我们称之为可预测过程。
最简单的此类过程就是时间本身，即 \(x(t) = t\)。

更一般地，如果一个过程 \(\vx(t)\) 在每个时刻 \(t\)，其在 \(t + \delta t\) 的值
由其在 \(t\) 的值决定（其中 \(\delta t\) 是一个无穷小的增量），那么该过程就是\textit{可预测的}（predictable）。
通常，\(\vx(t)\) 是连续且光滑的，因此变化量
\(\delta \vx(t) = \vx(t + \delta t) - \vx(t)\) 是无穷小的。
此类过程通常由（多元）微分方程描述：
\begin{equation}
    \dot{\vx}(t) = f(\vx(t)), \quad \vx \in \R^{d}.
    \label{eqn:process}
\end{equation}
在系统理论 \cite{Cal:Des,Sastry-Nonlinear} 中，方程
\eqref{eqn:process} 被称为状态空间模型。一个受控过程由下式给出：
\begin{equation}
    \dot{\vx}(t) = f(\vx(t), \vu(t)), \quad \vx \in \R^{d}, \vu \in \R^{k},
    \label{eqn:process-controlled}
\end{equation}
其中 \(\vu(t)\) 是一个可计算的输入过程。

\begin{example}
    牛顿第二定律预测了运动物体在受力 \(\vF(t) \in \R^{3}\) 作用下的轨迹 \(\vx(t) \in \R^{3}\)：
    \begin{equation}
        m\ddot{\vx}(t) = \vF(t).
        \label{eqn:newton-second}
    \end{equation}
    当没有受力时，即 \(\vF(t) \equiv \vzero\)，这
    退化为牛顿第一定律：物体以恒定速度 \(\vv \in \R^{3}\) 运动：
    \begin{equation}
        \ddot{\vx}(t) = \vzero \iff \ \text{there exists } \vv \in
        \R^{3} \text{ such that } \dot{\vx}(t) = \vv.
        \label{eqn:newton-first}
    \end{equation}
\end{example}

\subsection{低维性}\label{sec:intro-low-dimensionality}
\paragraph{学习预测。}
现在假设你已经观察到或被给定了许多序列片段：
\begin{equation}
    \{S_{1}, S_{2}, \ldots, S_{i}, \ldots, S_{N}\}
\end{equation}
这些片段提取自一个可预测序列 \(\{x_{n}\}_{n=1}^{\infty}\)。
不失一般性，假设每个片段的长度为 \(D \gg d\)，因此
\begin{equation}
    S_{i} = [x_{j(i)}, x_{j(i)+1}, \ldots, x_{j(i)+D-1}]^{\top} \in \R^{D}
\end{equation}
其中 \(j \in \N\)。然后给你一个新的片段 \(S_{t}\)，并
要求你预测其未来的值。

困难在于生成函数 \(f\) 及其阶数 \(d\) 是未知的：
\begin{equation}
    x_{n+d} = f(x_{n+d-1}, \ldots, x_{n}).
    \label{eqn:sequence-order-d}
\end{equation}
因此，目标是从样本片段 \(S_{1}, S_{2}, \ldots, S_{N}\) 中学习 \(f\) 和 \(d\)。
学习预测的核心任务是：
\begin{center}
    \textit{给定一个可预测序列的许多采样片段，
    我们如何才能有效且高效地识别出函数 \(f\)？}
\end{center}

\paragraph{可预测性与低维性。}
为了识别预测函数 \(f\)，我们可能会注意到由 \eqref{eqn:sequence-order-d} 给出的任何可预测序列的片段都有一个共同特征。
如果我们取一个长片段，比如长度为 \(D \gg d\)，并将其视为一个向量
\begin{equation}
    \vx_{i} = [x_{i}, x_{i+1}, \ldots, x_{i+D-1}]^{\top} \in \R^{D},
    \label{eqn:segment-vector}
\end{equation}
那么所有此类向量 \(\{\vx_{i}\}\) 的集合绝不是随机的，
也不可能占据整个空间 \(\R^{D}\)。相反，它最多具有
\(d\) 个自由度——给定任何 \(\vx_{i}\) 的前 \(d\) 个元素，
其余元素就被唯一确定了。换句话说，
所有的 \(\{\vx_{i}\}\) 都位于一个 \(d\) 维曲面上。在
数学中，这样的曲面被称为子流形（submanifold），记为 \(\cS
\subset \R^{D}\)。

\begin{figure}[t]
    \centering
    \includegraphics[width=0.6\linewidth]{\toplevelprefix/chapters/chapter1/figs/two-dimensional
    plane in R10.pdf}
    \caption{十维环境空间中的二维子空间。}
    \label{fig:lowdimplane}
\end{figure}
在实践中，如果我们选择的片段长度 \(D\) 足够大，所有
从同一个预测函数中采样的片段都会位于一个本征维度为 \(d\) 的曲面上，
该维度显著低于环境空间的维度 \(D\)。
例如，如果序列由线性自回归给出：
\begin{equation}
    x_{n+2} = a x_{n+1} + b x_{n},
    \label{eqn:sequence-2d}
\end{equation}
其中 \(a, b \in \R\) 为常数，并且我们采样长度为 \(D = 10\) 的片段，
那么所有样本都位于 \(\R^{10}\) 中的一个二维平面上，
如\Cref{fig:lowdimplane}所示。识别出
这个二维子空间就完全确定了 \eqref{eqn:sequence-2d} 中的常数 \(a\)
和 \(b\)。

更一般地，当预测函数 \(f\) 是线性的时，正如在
\eqref{eqn:lineary-system} 和 \eqref{eqn:lineary-system-closed} 给出的系统中那样，
长片段总是位于一个低维线性子空间上。
此时，识别预测函数在很大程度上等同于识别这个子空间，
这个问题被称为主成分分析（principal component analysis）。
我们将在\Cref{ch:classic-models}中讨论这些经典模型和方法。

这一观察结果可以推广到一般的可预测序列：如果我们能够
识别出片段样本所在的低维曲面，
我们就能识别出预测函数 \(f\)。\footnote{在
    温和的条件下，低维曲面与函数 \(f\) 之间存在一一映射。这一事实
    已被应用于系统辨识等问题中，我们将在
稍后讨论。}我们再怎么强调这一性质的重要性都不为过：\textit{可预测序列的
长片段的所有样本都位于一个低维子流形上。}在很大程度上，
现代科学和数学正是试图识别和
描述观测现象的这种低维结构。这
解释了为什么关于我们物理世界的最重要和最著名的科学发现
通常用方程来描述，如图
\ref{fig:famous-equations} 所示，因为方程将原本独立的观测物理量限制在
低维流形上。
\begin{figure}
    \centering
    \includegraphics[width=0.65\linewidth]{\toplevelprefix/chapters/chapter1/figs/Equations.png}
    \caption{改变世界的17个著名方程。其中大多数
        方程描述了宇宙的物理定律。对于
        智能而言，第15个关于熵的方程将扮演非常特殊的
        角色，我们将在本书中解释原因，特别是在第
    \ref{ch:denoising} 章中。}
    \label{fig:famous-equations}
\end{figure}

根据统一了我们所知几乎所有物理定律的弦理论，
整个宇宙在一个10维空间中演化，
始于大爆炸，见图
\ref{fig:universe-10D}。正如我们将在
本书中看到的，所有现代学习方法也都隐式或显式地利用了
数据中的低维性。
\begin{figure}
    \centering
    \includegraphics[width=0.65\linewidth]{\toplevelprefix/chapters/chapter1/figs/universe-10D.jpg}
    \caption{宇宙的演化，基于迄今为止最好的
    物理理论，被认为发生在一个10维空间中。}
    \label{fig:universe-10D}
\end{figure}

\begin{remark}[物理学中的熵与智能中的熵]
    在图
    \ref{fig:famous-equations} 中，我们看到唯一一个{\em 没有}
    以方程形式描述的物理定律是
    热力学第二定律，它指出任何封闭物理系统（如整个宇宙）的熵总是{\em
    增加}的。我们将在本书中看到，智能（和科学）的目标是
    在观测现象中识别出能够{\em 降低}我们预测中熵的结构。
    这将在第
    \ref{ch:denoising} 章和附录 \ref{app:entropy} 中进行相当广泛的研究。
\end{remark}

在现实世界的场景中，观测数据通常来自多个
可预测序列。例如，一个视频序列可能包含
几个运动的物体。在这种情况下，数据位于
低维线性子空间或非线性子流形的混合体上，如
\Cref{fig:mixture-models} 所示。
\begin{figure}
    \centering
    \includegraphics[width=0.8\linewidth]{\toplevelprefix/chapters/chapter1/figs/mixture.pdf}
    \caption{分布在（正交）子空间（左）或子流形（右）混合体上的数据。}
    \label{fig:mixture-models}
\end{figure}

\paragraph{低维性的性质。}
可预测序列中的时间相关性并不是数据呈现低维性的唯一原因。
例如，所有图像构成的空间是
巨大的，但其中大部分由无结构的随机图像组成，如
\Cref{fig:noise-image}（左）所示。然而，自然图像和视频
由于像素值之间存在强烈的空间和时间相关性，因此具有高度的冗余性。
这种冗余性使我们能够
轻松识别图像是嘈杂的还是干净的，如
\Cref{fig:noise-image}（中和右）所示。因此，
相对于像素总数，自然图像的分布具有非常低的本征维度。

\begin{figure}
    \centering
    \includegraphics[height=0.30\linewidth]{\toplevelprefix/chapters/chapter1/figs/Gaussian-noise.png}\hspace{2mm}
    \includegraphics[height=0.30\linewidth]{\toplevelprefix/chapters/chapter1/figs/Standard-test-image-Barbara-of-size-512-512-pixels-including-Gaussian-noise-with.png}
    \hspace{2mm}
    \includegraphics[height=0.30\linewidth]{\toplevelprefix/chapters/chapter1/figs/barbara.jpg}
    \caption{随机噪声图像（左）与含噪图像（中）及原始干净图像（右）的对比。}
    \label{fig:noise-image}
\end{figure}

由于学习低维结构既重要又
普遍，因此《基于低维模型的高维数据分析：原理、计算与应用》（\textit{High-Dimensional Data Analysis with
Low-Dimensional Models: Principles, Computation, and Applications}）
\cite{Wright-Ma-2022} 一书开篇即指出：“\textit{在
    高维空间中识别信号或数据的低维结构问题是最基本的问题之一，
    在漫长的历史中，它交织了系统理论、信号处理、
模式识别、机器学习和统计学等众多工程和数学领域。}”

通过将观测数据点 $\vx$ 约束在
低维曲面上，我们使其各个元素高度依赖或
相互关联，并且在某种意义上，可以从
其他元素的值中“预测”出来。例如，如果我们知道数据被
约束在 $\R^D$ 中的 $d$ 维曲面上，除了预测之外，我们还可以执行
几项有用的任务：
\begin{itemize}
    \item \textbf{补全（completion）}：给定一个典型样本 $\vx$ 的多于 $d$ 个元素，
        其余元素通常可以被
        唯一确定；\footnote{预测成为这一性质的一个特例。}
    \item \textbf{去噪（denoising）}：如果样本 $\vx$ 的元素受到
        小噪声的扰动，通过将 $\vx$ 投影回曲面上，可以有效地
        去除噪声；
    \item \textbf{纠错（error correction）}：如果 $\vx$ 的少量未知
        元素被任意损坏，它们可以被
        高效地纠正。
\end{itemize}
\Cref{fig:low-dim-properties} 使用一个
低维线性结构——二维平面中的一维直线——来说明这些性质。

\begin{figure}
    \centering
    \includegraphics[height=0.28\linewidth]{\toplevelprefix/chapters/chapter1/figs/Completion-low-dim.png}
    \includegraphics[height=0.28\linewidth]{\toplevelprefix/chapters/chapter1/figs/Denoising-low-dim.png}
    \includegraphics[height=0.32\linewidth]{\toplevelprefix/chapters/chapter1/figs/Correction-low-dim.png}
    \caption{低维（线性）结构性质的图示：
        它能够实现补全（左）、去噪（中）和
    纠错（右）。}
    \label{fig:low-dim-properties}
\end{figure}

在温和的条件下，这些性质可以推广到高维空间中的许多其他
低维结构 \cite{Wright-Ma-2022}。正如我们将看到的，这些有用的
性质——例如补全和去噪——启发了
学习此类结构的有效方法。

为了简单起见，到目前为止，我们一直使用确定性情况来
介绍可预测性和低维性的概念，其中
数据精确地位于子空间或曲面等几何结构上。
然而，在实践中，数据总是包含一些不确定性
或随机性。在这种情况下，我们可以假设数据服从
密度为 \(p(\vx)\) 的概率分布。如果一个分布的密度集中在
一个低维几何结构（子空间、曲面或它们的混合体）周围，
如\Cref{fig:mixture-models}所示，那么该分布就被认为是“低维的”。一旦
学习到，这样的密度 \(p(\vx)\) 就可以作为强大的先验，用于
从部分、嘈杂或损坏的观测中估计 \(\vx\)：
\begin{equation}
    \vy = f(\vx) + \vn,
    \label{eqn:noisy-observation}
\end{equation}
通过计算条件估计 \(\hat{\vx}(\vy) = \Ex(\vx \mid
\vy)\) 或通过对条件分布 \(\hat{\vx}(\vy)
\sim p(\vx\mid \vy)\) 进行采样。\footnote{现代生成式 AI 技术，
    如条件图像生成，严重依赖于这一原理，
正如我们将在\Cref{ch:conditional-inference}中讨论的那样。}

这些讨论引出了本书
理解智能和深度网络的演绎方法所基于的唯一假设：
\begin{quote}
    \textbf{主要假设：}\textit{任何智能系统或
        学习方法都应该（并且能够）依赖于世界的
        可预测性；因此，观测到的高维数据样本的分布
    具有低维支撑。}
\end{quote}
剩下的问题是如何从高维数据中正确且计算高效地学习这些低维结构。
正如我们将看到的，经典分析方法中研究的参数化模型
与现代实践中流行的非参数化模型（如深度网络），
仅仅是殊途同归的不同手段。

\section{如何学习？}

\subsection{基于分析模型的方法}
\label{sec:analytical}
请注意，即使一个预测函数在计算上是易于处理的，也
并不意味着从有限数量的采样片段中学习这个函数是易于处理或可扩展的。
确保易处理性的一种经典方法是，对我们正在处理的低维结构族做出明确的假设。
在历史上，由于计算能力和数据的限制，简单且
理想化的分析模型最先被研究，因为它们
通常提供高效的闭式解或数值解。此外，
它们为更一般的问题提供了深刻的见解，并且
已经为重要（尽管有限）的情况产生了有用的解决方案。
在过去计算资源匮乏的时代，允许高效闭式解或数值解的分析模型
是唯一能够被实现的方案。\textit{线性
结构}成为第一类被彻底研究的模型。

例如，可以说最简单的情况是假设数据
分布在高维空间中的单个低维子空间周围。
在某种程度上等价地，人们可以假设
数据服从一个几乎退化的低维高斯分布。
从有限数量的（含噪）样本中识别这样的子空间或高斯分布
是主成分分析（PCA）的经典问题，
并且已经为这类模型开发了有效的算法 \cite{JolliffeI2002}。人们可以
使模型族变得越来越复杂和富有表现力。
例如，人们可以假设数据分布在
低维分量（子空间或低维高斯分布）的混合体周围，
如独立成分分析（ICA）
\cite{Ans-1985}、字典学习（DL）、广义主成分
分析（GPCA）\cite{Vidal-GPCA}，或者近年来在压缩感知等领域被广泛研究的更一般的
稀疏低维模型类
\cite{Wright-Ma-2022}。

在所有这些分析模型族中，核心问题是
在每个族中识别出最适合给定数据的最紧凑模型。
下面，我们简要介绍这些经典的
分析模型，但将更系统的研究留到
\Cref{ch:classic-models}。在理论上，这些分析模型为
低维结构的几何和统计性质提供了
巨大的洞察力。它们通常产生闭式解或
高效且可扩展的算法，这对于其分布可以被此类模型很好地近似的数据
非常有用。更
重要的是，对于更一般的问题，它们提供了一种直觉，即识别低维结构的问题
可能有多容易或多困难，以及解决此类问题的基本思想是什么。

\subsubsection{线性动力系统} \label{sec:linear-systems}

\vspace{2mm}
\paragraph{维纳滤波器（Wiener filter）。}
正如在\Cref{sec:predictability}中所讨论的，智能的一个核心任务是
学习观测序列中什么是可预测的。
最简单的一类可预测序列——或
信号——是由\textit{线性时不变}（LTI）过程生成的：
\begin{equation}
    x[n] = h[n]*z[n] + \epsilon[n],
    \label{eqn:Wiener-model}
\end{equation}
其中 \(*\) 是卷积操作，\(z\) 是输入，
\(h\) 是脉冲响应函数。\footnote{通常，假设 \(h\) 具有
    某些理想的性质，例如有限长度
或带限频谱。} 这里 \(\epsilon[n]\) 表示加性
观测噪声。给定输入过程 \(\{z[n]\}\) 和
输出过程的观测值 \(\{x[n]\}\)，目标是找到
最优的 \(h[n]\)，使得 \(\hat{x}[n] = h[n]*z[n]\) 能够最优地预测
\(x[n]\)。预测质量（即优度）由
均方误差（MSE）来衡量：
\begin{equation}
    \min_{h} \Ex[\epsilon[n]^{2}] = \Ex[\|x[n] - h[n]*z[n]\|_{2}^{2}].
    \label{eqn:wiener-mse}
\end{equation}
最优解 \(h[n]\) 被称为（去噪）滤波器。诺伯特·维纳（Norbert
Wiener）——他也是控制论运动的创始人——在20世纪40年代研究了这个问题，并推导出了一个优雅的闭式解，
被称为\textit{维纳滤波器} \cite{Wiener-1942,Wiener-1949}。
这一结果，也被称为最小方差估计和滤波，
成为了信号处理的基石之一。感兴趣的
读者可以参考 \cite[附录 B]{MaY2003} 以获取此类估计器的详细推导。

\paragraph{卡尔曼滤波器（Kalman filter）。}
对动力系统进行去噪或滤波的思想后来在20世纪60年代被鲁道夫·卡尔曼（Rudolph Kalman）扩展到了由有限维状态空间模型描述的线性时不变系统：
\begin{equation}
    \vz[n] = \vA \vz[n-1] + \vB\vu[n] + \vepsilon[n].
    \label{eqn:linear-state-space}
\end{equation}
问题是从以下形式的含噪观测中估计系统状态 \(\vz[n]\)：
\begin{equation}
    \vx[n] = \vC \vz[n] + \vw[n],
    \label{eqn:Kalman-model}
\end{equation}
其中 \(\vw\) 是（白）噪声。最优的因果\footnote{这
    意味着估计器只能使用直到当前时间步 \(n\) 的观测值。卡尔曼滤波器总是因果的，而维纳
滤波器则不一定。}状态估计器，
它最小化最小均方误差
（MMSE）预测误差
\begin{equation}
    \min \Ex[ \|\vx[n] - \vC\vz[n]\|_{2}^{2}]
    \label{eqn:kalman-mse}
\end{equation}
由所谓的\textit{卡尔曼滤波器} \cite{kalman1960new} 以闭式解给出。
这一结果是现代控制理论的基石，
因为它使得从含噪观测中估计动力系统的状态成为可能。
然后，人们可以引入线性状态反馈，
例如 \(\vu[n] = \vF \hat{\vz}[n]\)，并使
闭环系统完全自治，如
\Cref{eqn:recursive-closed-loop} 所示。感兴趣的读者可以参考
\cite[附录 B]{MaY2003} 以获取卡尔曼滤波器的详细推导。

\paragraph{线性动力系统的辨识。}
为了推导卡尔曼滤波器，假设系统参数 \(\vA, \vB, \vC\)
是已知的。如果它们没有预先给定，
问题就会变得更具挑战性，这被称为\textit{系统
辨识}（system identification）：如何从输入序列 \(\{\vu[n]\}\) 和
观测序列 \(\{\vx[n]\}\) 的许多样本中\textit{学习}参数 \(\vA, \vB,
\vC\)。这是系统理论中的一个经典问题。
如果系统是线性的，输入和输出
序列 \(\{\vu[n], \vx[n]\}\) 共同位于某个
低维子空间上\footnote{该子空间的维度与
状态空间模型 \eqref{eqn:linear-state-space} 的阶数相同。}。
因此，辨识问题本质上等同于
识别这个低维子空间
\cite{OverscheeP1996,Liu-2009-CDC,Liu-2010-SIAM}。

请注意，上述问题有两个共同点：首先，
无噪序列或信号被假设是由一个
显式的参数化模型族生成的；其次，这些模型
本质上是线性的。从概念上讲，设 \(\vx_{o}\) 是一个随机变量，
其“真实”分布支撑在一个
低维线性子空间（比如 \(S\)）上。在很大程度上，
维纳滤波器和卡尔曼滤波器都试图从其含噪观测中估计这样一个
\(\vx_{o}\)：
\begin{equation}
    \vx = \vx_{o} + \vepsilon, \quad \vx_{o} \sim S,
\end{equation}
其中 \(\vepsilon\) 通常是随机高斯噪声（或
过程）。因此，它们的解决方案依赖于识别一个
最适合观测到的含噪数据的低维线性子空间。
通过将数据投影到这个子空间上，人们可以获得
最优的去噪操作，所有这些都以闭式解的形式呈现。

\subsubsection{线性与混合线性模型}
\label{sec:PCA-ICA}
\paragraph{主成分分析（Principal component analysis）。}
从上述经典信号处理和系统辨识的问题中，
我们可以看到，所有这些问题背后的主要任务是
从含噪观测中学习一个\textit{单一的}低维线性子空间。
在数学上，我们可以将这种结构建模为
\begin{equation}
    \vx = \vu_{1} z_{1} + \vu_{2} z_{2} + \cdots + \vu_{d} z_{d} +
    \vepsilon =  \vU \vz + \vepsilon, \quad \vU \in \R^{D\times d},
    \label{eqn:PCA-model}
\end{equation}
其中 \(\vepsilon \in \R^{D}\) 是小的随机噪声。\Cref{fig:PCA}
展示了具有两个主成分的这种分布。
\begin{figure}
    \centering
    \includegraphics[width=0.5\linewidth]{\toplevelprefix/chapters/chapter1/figs/PCA.png}
    \caption{具有两个主成分的分布。}
    \label{fig:PCA}
\end{figure}
问题是从 \(\vx\) 的许多样本中找到子空间基 \(\vU\)。
一种典型的方法是最小化到子空间上的投影误差：
\begin{equation}
    \min_{\vU}\Ex[\|\vx - \vU\vU^{\top} \vz\|_{2}^{2}].
\end{equation}
这本质上是一个去噪任务：一旦正确找到了基 \(\vU\)，
我们就可以通过将含噪样本 \(\vx\) 投影到由 \(\vU\) 张成的低维子空间上来对其进行去噪：
\begin{equation}
    \vx \rightarrow \hat{\vx} = \vU\vU^{\top} \vx.
\end{equation}
如果噪声很小且低维子空间 \(\vU\) 被正确学习，我们期望 \(\vx \approx \hat{\vx}\)。因此，PCA
是所谓的“自编码”（auto-encoding）方案的一个特例，该方案通过将数据投影到低维空间来对其进行编码，并通过将其投影回原始空间来对低维代码进行解码：
\begin{equation}
    \vx   \xrightarrow{\hspace{2mm} \vU^{\top}\hspace{2mm}} \vz
    \xrightarrow{\hspace{2mm} \vU \hspace{2mm}} \hat{\vx}.
    \label{eqn:auto-encoding-PCA}
\end{equation}
由于数据结构简单，编码器 \(\cE\) 和解码器 \(\cD\) 都变成了简单的线性算子（投影和提升）。

统计学中的这个经典问题被称为主成分分析（PCA）。它最早由 Pearson 在 1901 年研究
\cite{Pearson1901}，后来由 Hotelling 在 1933 年独立研究
\cite{Hotelling1933}。经典著作 \cite{Jolliffe1986,JolliffeI2002} 对该主题进行了系统的总结。
人们也可以显式地假设数据 \(\vx\) 服从单一的低维高斯分布：
\begin{equation}
    \vx \sim \mathcal{N}(\vzero, \vU\vU^{\top} + \sigma \vI), \quad
    \vU \in \R^{D\times d},
\end{equation}
这等价于假设 PCA 模型 \eqref{eqn:PCA-model} 中的 \(\vz\) 服从标准正态分布。这个问题被称为概率 PCA
\cite{TippingM1999}，并且具有与 PCA 相同的计算解。

在本书中，我们在 \Cref{ch:classic-models} 中从学习低维分布的角度重新审视
PCA。我们的目标是利用这个简单且理想化的模型来传达学习低维分布紧凑表征的一些最基本的思想，包括通过去噪进行压缩以及为了获得一致表征而进行自编码的重要概念。
\paragraph{独立成分分析。}
独立成分分析（ICA）最初由 \cite{Ans-1985}
提出，作为\textit{学习良好表征}的经典模型。事实上，它最初是作为我们记忆的一个简单数学模型而提出的。ICA 模型采取了与上述
PCA 模型 \eqref{eqn:PCA-model} 极具迷惑性的相似形式，它假设观测到的随机变量 \(\vx\) 是多个独立成分
\(z_{i}\) 的线性叠加：\footnote{在 PCA 中，“成分”是学习到的向量，而在 ICA 中它们是标量。这仅仅是名称和约定上的差异。}
\begin{equation}
    \vx = \vu_{1} z_{1} + \vu_{2} z_{2} + \cdots + \vu_{d} z_{d}  +
    \vepsilon =  \vU \vz + \vepsilon.
    \label{eqn:ICA-model}
\end{equation}
然而，这里的成分 \(z_{i}\) 被假设为独立的\textit{非高斯}变量。例如，一个常见的选择是
\begin{equation}
    z_{i} = \sigma_{i} \cdot w_{i}, \quad \sigma_{i} \sim \dBern(p),
    \label{eqn:ICA-modes}
\end{equation}
其中 \(\sigma_{i}\) 是伯努利随机变量，\(w_{i}\)
可以是常数值或另一个随机变量（例如高斯变量）。\footnote{即使 \(w\) 是高斯变量，\(\sigma w\)
也不再是高斯变量！} ICA 问题旨在从观测到的 \(\vx\) 样本中识别出 \(\vU\) 和 \(\vz\)。
\Cref{fig:ICA-PCA} 说明了 ICA 和 PCA 之间的区别。

\begin{figure}
    \centering
    \includegraphics[width=0.7\linewidth]{\toplevelprefix/chapters/chapter1/figs/PCA_ICA.png}
    \caption{PCA（左）与 ICA（右）的对比。请注意，PCA 寻找近似张成数据的\textit{正交}向量，而 ICA
    寻找\textit{独立}组合以近似形成数据的（不一定正交的）向量，以及这种组合的系数。}
    \label{fig:ICA-PCA}
\end{figure}

尽管一旦学习到 \(\vU\) 和 \(\vz\)，从 \(\vz\) 到 \(\vx\) 的（解码）映射似乎是线性且直接的，但从
\(\vx\) 到 \(\vz\) 的（编码）映射可能会很复杂，并且可能无法用简单的线性映射来表示。因此，ICA 通常给出以下形式的自编码：
\begin{equation}
    \vx   \xrightarrow{\hspace{2mm} \cE\hspace{2mm}} \vz
    \xrightarrow{\hspace{2mm} \vU \hspace{2mm}} \hat{\vx}.
    \label{eqn:auto-encoding-ICA}
\end{equation}
因此，与 PCA 不同，ICA 的分析和求解要困难一些。在 20 世纪 90 年代，Erkki Oja 和 Aapo
Hyv\"{a}rinen 等研究人员 \cite{hyvarinen-1997,Hyvrinen-2000} 对 ICA
做出了重大的理论和算法贡献。在 \Cref{ch:classic-models} 中，我们将研究并提供 ICA 的一种解法，从中编码映射
\(\cE\) 将变得清晰。事实上，Aapo Hyv\"{a}rinen 在 2005 年提出了通用的得分匹配（score
matching）方法 \cite{hyvarinen05a}，部分原因是为了学习像 ICA 中那样的高斯混合分布。我们将在稍后以及第
\ref{ch:denoising} 章中对此进行更详细的讨论。

\paragraph{稀疏结构与压缩感知。}
如果 \eqref{eqn:ICA-modes} 中的 \(p\) 非常小，那么任何成分非零的概率都很小。在这种情况下，我们称
\(\vx\) 是稀疏生成的，并集中在维度为 \(k = p \cdot d\) 的线性子空间的并集上。因此，我们可以将 ICA
模型扩展到一类更一般的低维结构，即稀疏模型。

一个 \(k\)-稀疏模型是所有\textit{\(k\)-稀疏向量}的集合：
\begin{equation}
    \cZ = \{\vz \in \R^{n} \mid \|\vz\|_{0} \leq k\},
\end{equation}
其中 \(\|\cdot\|_{0}\) 计算非零元素的个数。因此 \(\cZ\) 是所有与坐标轴对齐的 \(k\) 维子空间的并集，如
\Cref{fig:mixture-models}（左）所示。信号处理和统计学中的一个核心问题是从其线性观测中恢复稀疏向量 \(\vz\)：
\begin{equation}
    \vx = \vA \vz + \vepsilon, \quad \vA \in \R^{m\times n},
    \label{eqn:sparse-model}
\end{equation}
其中 \(\vA\) 是给定的，通常 \(m < n\)，且 \(\vepsilon\) 是小噪声。这个看似无害的问题在求解上是 NP
难的，甚至难以近似（详见 \cite{Wright-Ma-2022}）。

尽管有着可追溯到 18 世纪的丰富历史 \cite{Boscovichca1750}，但对于此类问题并不存在可证明的高效算法，尽管在 20
世纪 60 年代到 90 年代之间提出了许多启发式算法。其中一些在实践中有效，但缺乏严格的理论证明。2000
年代初出现了一项重大突破，当时包括 David Donoho、Emmanuel Cand\'{e}s 和 Terence Tao
在内的数学家 \cite{donoho2005neighborly,Candes2005,CandesE2005-IT}
建立了一个严格的理论框架，刻画了通过凸 \(\ell^{1}\) 最小化有效且高效地解决稀疏恢复问题的精确条件：
\begin{equation}
    \min \|\vz\|_{1} \quad \text{subject to} \quad \|\vx -
    \vA\vz\|_{2} \leq \epsilon,
\end{equation}
其中 \(\|\cdot\|_{1}\) 是促进稀疏性的向量 \(\ell^{1}\) 范数，\(\epsilon\)
是一个小的常数。任何解都会产生一个稀疏（自）编码：
\begin{equation}
    \vx \xrightarrow{\hspace{2mm} \cE \hspace{2mm}} \vz
    \xrightarrow{\hspace{2mm} \vA \hspace{2mm}} \hat{\vx}.
    \label{eqn:decoding-sparse}
\end{equation}
我们将在 \Cref{ch:classic-models}
中描述这样一种算法（以及由此产生的映射），揭示稀疏编码与深度学习之间的基本联系。\footnote{早在 2010
年，人们就注意到了稀疏编码算法与深度网络之间的相似性 \cite{gregor2010learning}。}

\(\ell^{1}\) 最小化成功的条件出人意料地宽泛：成功恢复所需的最小测量次数 \(m\) 仅与内在维度 \(k\)
成正比。这就是\textit{压缩感知}（compressive sensing）现象
\cite{CandesE2006-ICM}。它扩展到了广泛的低维结构家族，包括低秩矩阵。这些结果从根本上改变了我们对在高维空间中恢复低维结构的理解。David
Donoho 等人将这种反转赞誉为“维度的祝福”（blessing of
dimensionality）\cite{DonohoD2000}，这与通常对“维度灾难”（curse of
dimensionality）的悲观信念形成了鲜明对比。完整的理论和结果体系在 \cite{Wright-Ma-2022} 中有详细介绍。

这一框架的计算意义怎么强调都不为过。它将以前被认为难以处理的问题转化为不仅易于处理，而且可以使用极其高效的算法进行可扩展求解的问题：
\begin{equation}
    \text{intractable} \;\Longrightarrow\; \text{tractable}
    \;\Longrightarrow\; \text{scalable}.
\end{equation}
在给定数据和计算的精确要求下，这些算法具有严格的正确性理论保证。这种方法的演绎性质与深度神经网络的经验性、归纳性实践形成了鲜明对比。然而，我们现在知道，这两种方法都有一个共同的目标——在高维空间中追求低维结构。

\paragraph{字典学习。}
从概念上讲，当观测矩阵 \(\vA\) 未知且必须从一组（可能带有噪声的）观测值 \(\vX = [\vx_{1}, \vx_{2},
\ldots, \vx_{n}]\) 中学习时，就会出现一个比稀疏编码问题 \eqref{eqn:sparse-model} 更难的问题：
\begin{equation}
    \vX = \vA \vZ + \vE, \quad \vA \in \R^{m\times n}.
    \label{eqn:dictionary-learning}
\end{equation}
这里我们只给定了 \(\vX\)，而没有对应的 \(\vZ = [\vz_{1}, \vz_{2}, \ldots, \vz_{n}]\)
或噪声项 \(\vE = [\vepsilon_{1}, \vepsilon_{2}, \ldots,
\vepsilon_{n}]\)，除了已知 \(\vz_{i}\) 是稀疏的。这就是\textit{字典学习}（dictionary
learning）问题，它推广了前面讨论的 ICA 问题 \eqref{eqn:ICA-model}。换句话说，假设数据 \(\vX\)
的分布是标准稀疏分布 \(\vZ\) 在线性变换 \(\vA\) 下的像，我们希望同时学习 \(\vA\) 及其“逆”映射
\(\cE\)，从而获得一个自编码器：
\begin{equation}
    \vX   \xrightarrow{\hspace{2mm} \cE\hspace{2mm}}  \vZ
    \xrightarrow{\hspace{2mm} \vA \hspace{2mm}} \hat{\vX}.
    \label{eqn:decoding-DL}
\end{equation}

PCA、ICA
和字典学习都假设数据分布支撑在低维线性或分段线性结构上或其附近。每种方法都需要从噪声样本中学习这些结构的（全局或局部）基。在
\Cref{ch:classic-models}
中，我们将研究如何通过这些经典模型来识别此类低维结构。特别是，我们将看到所有这些低维（分段）线性模型都可以通过同一类被称为\textit{幂迭代}（power
iteration）\cite{Zhai-2020}
的快速算法进行高效学习。尽管这些线性或分段线性模型对于大多数真实世界数据来说有些理想化，但理解它们是迈向理解更一般低维分布的必不可少的第一步。

\subsubsection{一般分布}\label{sec:denoising-intro}

图像、视频和音频等真实世界数据的分布过于复杂，无法用上述有些理想化的线性模型或高斯过程来建模。我们通常无法\textit{先验地}知道它们是由哪个参数模型族生成的。从历史上看，人们曾多次尝试为这些数据开发分析模型。特别是，菲尔兹奖得主
David Mumford 在 20 世纪 90 年代花费了大量精力试图理解和建模自然图像的统计特性
\cite{Mumford1996TheSD}。他及其学生（包括 Song-Chun
Zhu）从统计物理学中汲取灵感和技术，提出了许多用于自然图像分布的统计和随机模型，例如随机场或随机过程
\cite{Zhu-Entropy-1997,Zhu1997LearningGP,Zhu1997Prior,Huang-Mumford,Mumford-1999,Lee-Mumford}。然而，这些分析模型在生成与自然图像高度相似的样本方面只取得了有限的成功。显然，对于像图像这样的真实世界数据，我们需要开发更强大、更统一的方法来追求它们更一般的低维结构。

在实践中，对于真实世界数据的一般分布，我们通常只有来自该分布的大量样本——即所谓的\textit{经验分布}（empirical
distribution）。在这种情况下，我们通常不能指望它们的低维结构有一个清晰的解析形式，也不能指望由此产生的去噪算子有解析形式。\footnote{这与
PCA、维纳滤波器（Wiener filter）和卡尔曼滤波器（Kalman filter）的情况不同。}
因此，我们需要为这些经验分布开发一种更通用的解决方案，不一定是闭式解，但至少是可高效计算的。如果我们做得正确，上述线性模型的解应该成为它们的特例。

\paragraph{去噪。} 在 20 世纪 50 年代，统计学家开始对从任意分布中提取的数据进行去噪的问题产生兴趣。设
\(\vx_{o}\) 是一个概率密度函数为 \(p_{o}(\cdot)\) 的随机变量。假设我们观测到 \(\vx_{o}\) 的一个带噪版本：
\begin{equation}
    \vx = \vx_{o} + \sigma \vg,
\end{equation}
其中 \(\vg \sim \dNorm(\vzero, \vI)\) 是标准高斯噪声，\(\sigma\) 是观测的噪声水平。设
\(p(\cdot)\) 为 \(\vx\) 的概率密度函数。\footnote{即 \(p(\vx) =
        \int_{-\infty}^{\infty} \phi_{\sigma}(\vx - \vz)p_{o}(\vz)
    \odif{\vz}\)，其中 \(\phi_{\sigma}\) 是高斯分布 \(\mathcal{N}(\vzero,
\sigma^{2} \vI)\) 的密度函数。} 值得注意的是，给定 \(\vx\) 时 \(\vx_{o}\)
的后验期望可以通过一个优雅的公式来计算，该公式被称为 Tweedie 公式
\cite{Robbins1956AnEB}：\footnote{Herbert Robbins 根据他们之间的私人通信，将该公式归功于
Maurice Kenneth Tweedie。}
\begin{equation}
    \hat{\vx}_{o} = \Ex[\vx_{o}\mid \vx] = \vx + \sigma^{2} \nabla \log p(\vx).
\end{equation}
从公式中可以看出，函数 \(\nabla \log p(\vx)\) 在这里对观测值 \(\vx\) 的去噪中扮演着非常特殊的角色。噪声
\(\vg\) 可以显式地估计为
\begin{equation}
    \hat{\vg} = \frac{\vx - \hat{\vx}_{o}}{\sigma} = -\sigma \nabla \log p(\vx),
\end{equation}
为此，我们只需要知道 \(\vx\) 的分布 \(p(\cdot)\)，而不需要知道 \(\vx_{o}\) 的真实分布
\(p_{o}(\cdot)\)。这个结果的一个重要推论是，如果我们向任何分布添加高斯噪声，只要我们能以某种方式获得函数 \(\nabla
\log p(\vx)\)，去噪过程就可以轻松完成。

因为这是一个如此重要且有用的结果，它在许多不同的背景和领域中被重新发现和使用。例如，在 Tweedie 公式
\cite{Robbins1956AnEB} 之后，它在几年后被 \cite{Miyasawa61}
重新发现，并在那里被称为“经验贝叶斯去噪”（empirical Bayesian denoising）。在 2000 年代初，函数
\(\nabla \log p(\vx)\) 在学习一般分布的背景下再次被重新发现，并被 Aapo Hyv\"{a}rinen
命名为“得分函数”（score function）\cite{hyvarinen05a}。它与（经验贝叶斯）去噪的联系很快被
\cite{Vincent2011} 认识到。Eero Simoncelli 的团队 \cite{Raphan10}
将其推广到了其他测量分布（超越高斯噪声）。今天，Tweedie 公式和去噪最直接的应用是通过迭代去噪进行图像生成
\cite{Kadkhodaie21a,ho2020denoising}。

\begin{figure}
    \centering
    \includegraphics[width=1\linewidth]{\toplevelprefix/chapters/chapter1/figs/Density-compress.png}
    \caption{左图为具有密度 \(p(\vx)\) 的分布的得分函数 \(\nabla \log p(\vx)\)
        的几何解释。由得分函数生成的算子将分布推向更高密度的区域。其目标是，通过某种紧凑性度量（例如熵或编码长度），使得到的分布更加“压缩”。最终，分布收敛到一个具有低维支撑的分布，如右图中的
    \(p^{*}(\vx)\) 所示。}
    \label{fig:score-function}
\end{figure}

\paragraph{熵最小化。} 得分函数具有直观的信息论和几何解释。在信息论中，\(-\log p(\vx)\) 对应于编码
\(\vx\) 所需的比特数\footnote{至少对于离散变量是这样，详见 \Cref{ch:denoising}。}。梯度
\(\nabla \log p(\vx)\) 指向更高概率密度的区域，如
\Cref{fig:score-function}（左）所示。沿这个方向移动会减少编码 \(\vx\) 所需的比特数。因此，算子
\(\nabla \log p(\vx)\) 推动分布向高密度区域“收缩”。形式上，可以证明（微分）熵
\begin{equation}
    H(\vx) = - \int p(\vw) \log p(\vw) \odif{\vw}
\end{equation}
在此操作下会减小（见 \Cref{ch:denoising} 和
\Cref{app:entropy}）。使用最优码本，得到的分布实现了更低的编码率，因此更加压缩。无限重复这个去噪过程会产生一个概率质量集中在较低维度支撑上的分布。例如，\Cref{fig:score-function}（左）中的分布
\(p(\vx)\) 收敛到 \(p^{*}(\vx)\)（右）：\footnote{严格来说，\(p^{*}(\vx)\)
    是一个广义函数：\(p^{*}(\vx) = p^{*}(\vx_{1})\delta(\vx-\vx_{1}) +
    p^{*}(\vx_{2}) \delta(\vx-\vx_{2})\)，其中 \(p^{*}(\vx_{1}) +
        p^{*}(\vx_{2}) =
1\)。}\footnote{请注意，在本节中，我们讨论的是迭代去噪和压缩高熵噪声分布，直到其收敛到低熵数据分布的过程。正如我们将在
    \Cref{ch:denoising} 中看到的，这个问题与学习数据分布的最优编码问题是对偶的，这也是一个有用的应用
\cite{relic2024lossy}。}
\begin{equation}
    H(\vx) = - \int p(\vw) \log p(\vw) \odif{\vw} \quad
    \xrightarrow{\hspace{1mm} \text{decreasing} \hspace{1mm}} \quad
    H^{*}(\vx) = - \int p^{*}(\vw) \log p^{*}(\vw) \odif{\vw}.
\end{equation}
随着分布收敛到 \(p^{*}(\vx)\)，由于连续熵和离散熵定义之间的技术差异，其微分熵趋于负无穷大。
\Cref{ch:denoising} 使用统一的\textit{率失真}（rate-distortion）度量解决了这个问题。

在本章稍后以及
\Cref{ch:denoising,ch:lossy-compression}
中，我们将探讨这个简单的去噪-压缩框架如何统一在高维空间中学习低维分布（包括自然图像分布）的强大方法。

\subsection{经验数据驱动方法}
在实践中，很难用上一节讨论的理想化线性、分段线性或其他分析模型来对重要的真实世界数据（如图像、声音和文本）进行建模。因此，从历史上看，人们提出了许多非参数模型和方法来建模或捕获数据的经验分布，然后使用这些“模型”进行推理。这些模型，特别是深度神经网络，通常从生物神经系统中汲取灵感，因为动物或人类的大脑处理此类数据的效率和效果都非常惊人。与上一节介绍的分析模型相比，这些数据驱动模型是最不被理解的，尽管在过去十年左右的时间里，它们在很大程度上主导了人工智能技术的实践。因此，本书的主要目标之一是提供一个有原则的统一框架（主要通过第
\ref{ch:deep-networks} 章到第 \ref{ch:applications}
章），以理解这些“现代”经验模型和经典分析模型，并揭示它们之间的基本联系。

\subsubsection{经典人工神经网络}
\paragraph{人工神经元。}

\begin{figure}[t]
    \centering
    \includegraphics[width=0.55\linewidth]{\toplevelprefix/chapters/chapter1/figs/neuron.png}
    \hspace{3mm}
    \includegraphics[width=0.40\linewidth]{\toplevelprefix/chapters/chapter1/figs/Artificial_neuron.png}
    \caption{第一个人工神经元的数学模型（右），它模拟了神经元（左）处理信号的方式。}
    \label{fig:neuron}
\end{figure}

受大脑神经系统的启发，第一个人工神经元\footnote{被称为线性阈值单元（Linear Threshold
Unit）或感知机（perceptron）。}的数学模型由 Warren
McCulloch\footnote{当时是芝加哥大学的精神病学教授} 和 Walter Pitts 于 1943 年提出
\cite{McCulloch-Pitts}。它将输入 $x_i$ 和输出 $o_j$ 之间的关系描述为：
\begin{equation}
    o_j = \varphi\Big( \sum_i w_{ji}x_i\Big),
\end{equation}
其中 $\varphi(\cdot)$ 是某种非线性激活函数，通常由阈值函数建模。该模型如 \Cref{fig:neuron}
所示。正如我们所见，这种形式已经具备了现代深度神经网络中基本单元的主要特征。该模型源于对我们神经系统中单个神经元工作方式的观察。然而，研究人员并不确切知道这些神经元的集合能够实现和执行什么功能。在更技术的层面上，他们也不确定应该使用哪种非线性激活函数
$\varphi(\cdot)$。因此，历史上提出了许多变体。\footnote{阶跃函数、硬阈值或软阈值、修正线性单元（ReLU）、sigmoid
等 \cite{dubey2022activation}。}

\begin{figure}[t]
    \centering
    \includegraphics[width=0.85\linewidth]{\toplevelprefix/chapters/chapter1/figs/single-deep.png}
    \caption{具有一个隐藏层的网络（左）与深度网络（右）的对比。}
    \label{fig:single-deep}
\end{figure}

\paragraph{人工神经网络。}
在 20 世纪 50 年代，Frank Rosenblatt 首次构建了一台包含此类人工神经元\textit{网络}的机器，如
\Cref{fig:perceptron} 所示。这台名为 Mark I Perceptron 的机器由一个输入层、一个输出层和一个包含
512 个人工神经元的单一隐藏层组成，如 \Cref{fig:perceptron} 左图所示，这与
\Cref{fig:single-deep}（左）中说明的类似。它被设计用于对字母的光学图像进行分类。然而，单层网络的容量是有限的，只能学习线性可分的模式。在
Marvin Minsky 和 Seymour Papert 1969
年出版的《感知机：计算几何简介》（\textit{Perceptrons: An Introduction to
Computational Geometry}）一书 \cite{Minsky-1969} 中，证明了 Mark I Perceptron
的单层架构无法学习 XOR（异或）函数。这一结果极大地削弱了人们对人工神经网络的兴趣，尽管后来证明多层网络可以学习 XOR 函数
\cite{Rumelhart1986}。事实上，如
\Cref{fig:single-deep}（右）所示，由这种简单神经元组成的足够大的多层网络可以模拟任何有限状态机，甚至是通用图灵机。\footnote{不要将神经网络在原则上能够做什么，与学习一个实现某些所需功能的神经网络是否易于处理或容易混为一谈。}
尽管如此，人工神经网络的研究随后在 20 世纪 70 年代进入了它的第一个寒冬。

\begin{figure}
    \centering
    \includegraphics[width=0.45\linewidth]{\toplevelprefix/chapters/chapter1/figs/visu-large.jpg}
    \hspace{2mm}
    \includegraphics[width=0.45\linewidth]{\toplevelprefix/chapters/chapter1/figs/Original-Mark-I-perceptron-as-seen-in-its-operators-manual-20.ppm.png}
    \caption{Frank Rosenblatt 在 20 世纪 50 年代末开发的 Mark I Perceptron 机器。}
    \label{fig:perceptron}
\end{figure}

\paragraph{卷积神经网络。}
20 世纪 50 年代和 60 年代对人工神经网络（如 Mark I
Perceptron）的早期实验有些令人失望。它们表明，像多层感知机（MLP）那样以一般方式简单地连接神经元可能是不够的。为了构建有效且高效的网络，理解网络中的神经元必须实现的集体目的或功能是非常有帮助的，这样它们就可以以专门的方式进行组织和学习。因此，在这个关头，机器智能的研究再次转向动物神经系统以寻求灵感。

众所周知，我们大脑的大部分都致力于处理视觉信息 \cite{PalmerS1999}。在 20 世纪 50 年代和 60 年代，David
Hubel 和 Torsten Wiesel
系统地研究了猫的视觉皮层。他们发现视觉皮层包含不同类型的细胞——简单细胞和复杂细胞——它们对不同方向和位置的视觉刺激很敏感
\cite{Hubel-Wiesel-1959}。Hubel 和 Wiesel 因这一突破性发现获得了 1981 年诺贝尔生理学或医学奖。

在人工神经网络方面，Hubel 和 Wiesel 的工作启发了 Kunihiko Fukushima（福岛邦彦）在 1980
年设计了“新认知机”（neocognitron）网络，该网络由模拟视觉皮层中生物神经元的人工神经元组成
\cite{Fukushima1980NeocognitronAS}。这被称为第一个\textit{卷积神经网络}（CNN），其架构如
\Cref{fig:neocognitron} 所示。与感知机不同，新认知机具有多个隐藏层，可以被视为一个深度网络，如
\Cref{fig:single-deep}（右）所示。

同样受猫视觉皮层中神经元工作方式的启发，Fukushima 首次引入了\textit{修正线性单元}（ReLU）：
\begin{equation}
    \varphi(x) = \max\{0, x\} = \casework{x, & \text{if } x > 0,
    \\ 0, & \text{if } x \leq 0,}
\end{equation}
作为 1969 年的激活函数 \(\varphi(\cdot)\)
\cite{Fukushima-1969}。然而，直到最近几年，ReLU
才在现代深度（卷积）神经网络中得到广泛应用。一旦我们讨论了深度网络实现的主要操作：压缩，本书将解释为什么 ReLU 是一个好的选择。

\begin{figure}
    \centering
    \includegraphics[width=0.6\linewidth]{\toplevelprefix/chapters/chapter1/figs/neocognitron.png}
    \caption{卷积神经网络的起源：
        由福岛邦彦（Kunihiko Fukushima）于1980年提出的神经认知机（Neocognitron）。请注意，
    卷积和池化的交替层模拟了在猫的视觉皮层中发现的简单细胞和复杂细胞的功能。}
    \label{fig:neocognitron}
\end{figure}

CNN类型的网络在20世纪80年代继续发展，许多不同的变体被提出并得到研究。然而，尽管深度网络具有非凡的能力，并且受神经科学启发改进了架构，但要针对图像分类等实际任务训练这样的深度网络仍然极其困难。让网络正常工作依赖于许多无法解释的启发式方法和技巧，这限制了神经网络的吸引力和适用性。大约在1989年，Yann
LeCun成功地使用\textit{反向传播}（back propagation,
BP）来学习用于识别手写数字的深度卷积神经网络\cite{LeCun-1989}，该网络后来被称为LeNet（见\Cref{fig:LeNet-5}），这标志着一项重大突破。经过几年的持续开发，他的坚持得到了回报：LeNet的性能最终在20世纪90年代末变得足够好，可以投入实际使用\cite{LeCun-1998}。它被美国邮政局用于识别手写数字（用于邮政编码）。LeNet被认为是所有现代深度神经网络（如我们稍后将讨论的AlexNet
\cite{krizhevsky2012imagenet}和ResNet
\cite{he2016deep}）的“原型”网络。因为这项工作，Yann
LeCun被授予了2018年图灵奖。\footnote{与另外两位深度网络的先驱Yoshua Bengio和Geoffrey Hinton共同获奖。}

\begin{figure}
    \centering
    \includegraphics[width=0.95\linewidth]{\toplevelprefix/chapters/chapter1/figs/LeNet-5.png}
    \caption{Yann LeCun在1989年设计的LeNet-5卷积神经网络。}
    \label{fig:LeNet-5}
\end{figure}

\paragraph{反向传播（Backpropagation）。}
纵观历史，深度神经网络的命运一直与其训练的难易程度和效率息息相关。反向传播（BP）正是为此而引入的。多层感知机可以表示为线性映射和非线性激活的复合：
\begin{equation}
    h(\vW_{1}, \ldots, \vW_{L}) = f^{L}(\vW_{L}f^{L-1}(\vW_{L-1}
    \cdots f^{2}(\vW_{2}f^{1}(\vW_{1}\vx)))).
\end{equation}
为了通过梯度下降训练网络权重\(\{\vW_{l}\}_{l=1}^{L}\)，我们必须计算梯度\(\partial h/\partial
\vW_{l}\)。微积分中的链式法则表明，可以高效地计算此类函数的梯度——这项技术后来被称为反向传播或BP；详情请参见\Cref{app:optimization}。在20世纪60年代和70年代，BP在最优控制和动态规划中就已为人所知，并出现在Paul
Werbos 1974年的博士论文中\cite{Werbos-1974, Werbos1994TheRO}。1986年，David
Rumelhart等人首次将其应用于训练多层感知机（MLP）\cite{Rumelhart1986}。从那时起，BP已成为训练深度网络的主导技术，因为它具有可扩展性，并且可以在并行和分布式计算平台上高效实现。然而，大自然很可能并不使用BP，\footnote{正如我们之前讨论的，大自然几乎无处不在地通过闭环反馈来学习纠正错误。}因为这种机制对于物理实现来说过于昂贵。\footnote{对于像大脑这样复杂的神经结构，端到端的BP在计算上是难以处理的，特别是如果更新需要实时发生的话。相反，对神经回路的一小部分进行局部更新在生物学上要可行得多。目前有相对少量的工作致力于将这种“局部学习”规则移植到训练深度网络中，从而规避BP
\cite{baldi2016theory,millidge2022predictive,li2025noprop}。我们认为这是提高训练深度网络可扩展性的一个令人兴奋的机会。}这为未来的改进留下了充足的空间。

尽管有上述算法上的进步，但在20世纪80年代和90年代，训练深度网络仍然非常繁琐且计算成本高昂。到了20世纪90年代末，支持向量机（SVM）\cite{SVM-1995}作为分类任务的卓越替代方案而广受欢迎。\footnote{分类问题的类似想法可以追溯到Thomas
    Cover
1964年的博士论文，该论文被浓缩并发表在1964年的一篇论文中\cite{Cover-1964}。}SVM提供了两个优势：被称为Vapnik--Chervonenkis（VC）理论的严格统计学习框架，以及高效的凸优化算法\cite{BoydVa04}。SVM的崛起在21世纪初迎来了神经网络的第二个寒冬。

\paragraph{压缩自编码（Compressive auto-encoding）。}
在20世纪80年代末和90年代，人工神经网络已经被用于学习图像等高维数据的低维表征。研究表明，神经网络可以直接从数据中学习PCA
\cite{Oja1982SimplifiedNM,Baldi89}，而不是使用\Cref{sec:PCA-ICA}中讨论的经典方法。在这一时期，人们还认为，由于神经网络能够对非线性变换进行建模，因此它们可以学习具有非线性分布的数据的低维表征。与线性PCA的情况类似，人们可以同时学习一个非线性降维编码器\(f\)和一个解码器\(g\)，它们各自都由深度神经网络建模\cite{Rumelhart1986,Kramer1991NonlinearPC}：
\begin{equation}
    \vX   \xrightarrow{\hspace{2mm} f \hspace{2mm}} \vZ
    \xrightarrow{\hspace{2mm} g \hspace{2mm}} \hat{\vX}.
    \label{eqn:auto-encoding-deep-networks}
\end{equation}
通过强制解码数据\(\hat{\vX}\)与原始数据\(\vX\)之间的一致性——例如，通过最小化MSE类型的重建误差\footnote{众所周知，MSE类型的误差对于具有复杂非线性结构的图像数据是有问题的。正如我们很快将要讨论的，包括GAN在内的生成方法的许多近期工作都集中在寻找原始数据\(\vX\)和重新生成的数据\(\hat{\vX}\)之间更好的替代距离函数。}：
\begin{equation}
    \min_{f,g} \|\vX - \hat{\vX}\|_{2}^{2}, \qquad \text{where} \quad
    \hat{\vX} = g(f(\vX)),
\end{equation}
可以直接从数据\(\vX\)中学习到一个自编码器。

但是，我们如何保证这种自编码确实捕获了\(\vX\)中真实的低维结构，而不是产生一个平凡的冗余表征呢？例如，人们可以简单地选择\(f\)和\(g\)为恒等映射，并设置\(\vZ
= \vX\)。为了确保自编码是有价值的，根据某种可计算的复杂性度量，所得到的表征应该是压缩的。1993年，Geoffrey
Hinton及其同事提出使用编码长度作为这样一种度量，将自编码的目标转化为寻找最小化编码长度的表征\cite{Hinton-1993}。这项工作在最小描述长度原则\cite{Rissanen-1978}和自由（亥姆霍兹）能最小化之间建立了根本的联系。Hinton团队后来的工作\cite{Hinton504}从经验上证明，这种自编码可以为真实世界的图像学习有意义的低维表征。Pierre
Baldi在2011年（就在深度网络获得广泛普及之前）提供了关于自编码器的全面综述\cite{Baldi2011}。我们将在\Cref{sec:unifying-approach}中进一步讨论复杂性度量和自编码，并在\Cref{ch:consistent-representations}中从一个更统一的视角对压缩自编码进行系统研究。

\begin{figure}
    \centering
    \includegraphics[width=0.8\linewidth]{\toplevelprefix/chapters/chapter1/figs/Comparison_image_neural_networks.svg.png}
    \caption{LeNet \cite{LeCun-1989}与AlexNet
    \cite{krizhevsky2012imagenet}的架构对比。}
    \label{fig:LeNet-AlexNet}
\end{figure}

\subsubsection{现代深度神经网络}
在近30年的时间里——从20世纪80年代到2010年代——神经网络并没有被主流机器学习社区认真对待。早期的深度网络（如LeNet）在数字识别等小规模分类问题上表现出了令人充满希望的性能，但它们的设计在很大程度上是经验性的，可用的数据集很小，而且反向传播对于那个时代的硬件来说计算成本高得令人望而却步。这些因素导致了人们兴趣的减退和进展的停滞，只有少数研究人员在坚持。

\paragraph{分类与识别。}
只有当充足的数据和计算能力变得可用时，深度网络的巨大潜力才能被释放出来。到了2010年代，诸如ImageNet之类的大型数据集已经出现，GPU也变得足够强大，使得即使对于比LeNet大得多的网络，反向传播的计算成本也变得可以承受。2012年，深度卷积网络AlexNet——以作者之一Alex
Krizhevsky的名字命名\cite{krizhevsky2012imagenet}——在ImageNet上以显著优势超越了现有的分类方法，从而引起了广泛关注。\footnote{深度网络已经在语音识别任务上取得了最先进的结果，但这些成功受到的关注远不如图像分类方面的突破。}\Cref{fig:LeNet-AlexNet}比较了AlexNet和LeNet。AlexNet保留了LeNet的许多特征——它只是规模更大，并将LeNet的sigmoid激活函数替换为ReLU。这项工作促成了Geoffrey
Hinton获得2018年图灵奖。

这一早期的成功激发了机器智能社区探索新的神经网络架构、变体和改进。从经验上发现，更大、更深的网络在图像分类等任务上能产生更好的性能。涌现出的著名架构包括VGG
\cite{Simonyan15}、GoogLeNet \cite{Szegedy2014GoingDW}、ResNet
\cite{He2016-lc}，以及最近的Transformers
\cite{vaswani2017attention}。尽管经验驱动的性能提升非常迅速，但对这些架构的理论解释——以及它们之间的关系（如果有的话）——仍然很匮乏。本书的一个目标是揭示这些网络优化的共同目标，并解释它们共有的特征，例如多层线性算子与非线性激活的交替（见\Cref{ch:deep-networks}）。

\paragraph{强化学习（Reinforcement learning）。}
早期深度网络的成功主要集中在语音和图像识别等监督分类任务上。后来，深度网络被采用来学习游戏博弈的决策或控制策略。在这种背景下，深度网络对最优决策/控制策略（即为了最大化预期奖励而采取的动作的概率分布）或相关的最优价值函数（对给定状态下预期奖励的估计）进行建模，如\Cref{fig:Alpha-Go}所示。网络参数根据使用当前策略玩游戏的成功或失败所返回的奖励进行增量优化\footnote{通常通过反向传播（BP）。}。这种学习范式被称为\textit{强化学习}
\cite{Sutton-Barto}；它起源于20世纪60年代末的控制系统实践\cite{Waltz1965AHA,Mendel1970ReinforcementlearningCA}，并可追溯到一段漫长而丰富的历史，包括Richard
Bellman的\textit{动态规划} \cite{Bellman-DP}和Marvin
Minsky在20世纪50年代的\textit{试错学习} \cite{Minsky-1954}。

从实现的角度来看，深度网络与强化学习的结合被证明是强大的：深度网络可以为难以进行解析建模的现实世界环境近似控制策略和价值函数。这在DeepMind的AlphaGo系统中达到了顶峰，该系统在2016年击败了顶尖围棋选手李世石，并在2017年击败了世界冠军柯洁，震惊了世界。\footnote{1996年，IBM的深蓝（Deep
    Blue）利用树搜索和剪枝等传统机器学习技术击败了俄罗斯国际象棋特级大师加里·卡斯帕罗夫（Garry
Kasparov），创造了历史，但这些方法已被证明可扩展性较差，并且在围棋等更复杂的游戏中并不成功。}

AlphaGo的成功令计算界感到惊讶，长期以来，计算界一直认为围棋的状态空间过于庞大，在计算和样本量方面无法找到任何高效的解决方案。唯一合理的解释是，围棋的最优价值和策略函数具有显著的有利结构：定性地说，它们的内在维度足够低，以至于可以被能够从可控数量的样本中学习的神经网络很好地近似。

\begin{figure}
    \centering
    \includegraphics[width=0.5\linewidth]{\toplevelprefix/chapters/chapter1/figs/Policy-Value.png}
    \caption{AlphaGo：使用深度神经网络对围棋游戏的最优策略或价值函数进行建模。}
    \label{fig:Alpha-Go}
\end{figure}

\paragraph{生成与预测。}
从某种角度来看，2010年代深度网络的早期实践主要集中在从数据\(\vX\)中提取相关信息，并将其编码为特定于任务的表征\(\vZ\)（假设\(\vZ\)表示分类任务中的类标签\footnote{这是一种非常不典型的标签符号——通常的符号是\(\vy\)——但出于我们的目的，将标签视为数据的高度压缩和稀疏表征是有用的。有关更多详细信息，请参见\Cref{example:image_classification_feature_learning}。}）：
\begin{equation}
    \vX   \xrightarrow{\hspace{2mm} f \hspace{2mm}} \vZ.
    \label{eqn:encoding-deep-networks}
\end{equation}
在这里，学习到的映射\(f\)不需要保留关于\(\vX\)的大部分分布信息；它只需保留任务的充分统计量即可。例如，样本\(\vx \in
\vX\)可能是一个苹果的图像，被\(f\)映射到类标签\(\vz = \text{``apple''}\)。

在许多现代设置中——例如所谓的大型通用（“基础”）模型——我们可能还需要解码\(\vZ\)以将相应的\(\vX\)恢复到规定的精度：
\begin{equation}
    \vZ   \xrightarrow{\hspace{2mm} g  \hspace{2mm}} \hat{\vX}.
    \label{eqn:decoding-deep-networks}
\end{equation}
因为\(\vX\)通常代表从外部世界观察到的数据，一个好的解码器允许我们模拟或预测世界上发生的事情。例如，在“文本到图像”或“文本到视频”任务中，\(\vz\)是描述所需图像\(\vx\)的文本，解码器应该生成一个内容与\(\vx\)匹配的\(\hat{\vx}\)。给定一个对象类别\(\vz
= \text{``apple''}\)，解码器\(g\)应该生成一个看起来像苹果的图像\(\hat{\vx}\)，尽管不一定与原始的\(\vx\)完全相同。

\paragraph{通过判别式方法进行生成。}
为了使生成的图像\(\hat{\vX}\)类似于真实的自然图像\(\vX\)，我们必须能够评估并最小化某种距离：
\begin{equation}
    \min \operatorname{dist}(\vX, \hat{\vX}).
\end{equation}
事实证明，对于具有低内在维度的高维空间中的分布，大多数具有理论动机的距离即使不是不可能，也是极难计算和优化的。\footnote{即使为\(\vX\)指定了参数化分布族，情况依然如此。对于具有低维支撑的分布，距离通常会变得病态或定义不明确。更糟糕的是，所选的分布族可能无法很好地近似我们感兴趣的真实分布。}

2007年，屠卓文（Zhuowen Tu）对早期试图对自然图像进行建模和生成的分析方法感到失望，决定尝试一种截然不同的方法。在CVPR
2007发表的一篇论文中\cite{Tu-2007}，他首次建议通过判别式方法学习图像的生成模型。这个想法很简单：如果评估距离\(\operatorname{dist}(\vX,
\hat{\vX})\)被证明很困难，那么可以转而学习一个判别器\(d\)来将\(\hat{\vX}\)与\(\vX\)区分开来：
\begin{figure}[h!]
    \centering
    \begin{tikzpicture}[node distance=1cm, auto]
        \node (Z) {$\Z$};
        \node (Xhat) [right=of Z] {$\hat{\X}$};
        \node (X) [below=of Xhat] {$\X$};
        \node (labels) [right=of Xhat] {$\{0, 1\}$};

        \draw[->] (Z) -- (Xhat) node[midway, above] {$g$};
        \draw[->] (Xhat) -- (labels) node[midway, above right] {$d$};
        \draw[->] (X) -- (labels) node[midway, below right] {$d$};
    \end{tikzpicture}
    \label{fig:gan-networks}
\end{figure}
其中标签\(0\)和\(1\)指示图像是生成的还是真实的。直观地说，区分\(\hat{\vX}\)和\(\vX\)变得越困难，它们可能就越接近。

屠卓文的工作\cite{Tu-2007}首次证明了从判别式方法学习生成模型的可行性。然而，该工作采用了传统的图像生成和分布分类方法（如boosting），这些方法被证明速度缓慢且难以实现。2012年之后，深度神经网络在图像分类领域受到欢迎。2014年，Ian
Goodfellow及其同事再次提出使用判别式方法生成自然图像\cite{Goodfellow-2014}。他们建议使用深度神经网络对生成器\(g\)和判别器\(d\)进行建模。此外，他们提出通过极小极大博弈（minimax
game）来学习\(g\)和\(d\)：
\begin{equation}
    \min_{g} \max_{d} \ell(\vX, \hat{\vX}),
\end{equation}
其中\(\ell(\cdot)\)是与分类任务相关的某种自然损失函数。在这个公式中，判别器\(d\)最大化其区分\(\vX\)和\(\hat{\vX}\)的成功率，而生成器\(g\)则做相反的事情。这种方法被称为\textit{生成对抗网络}（GANs）。研究表明，一旦在大型数据集上进行训练，GAN确实可以生成逼真的图像。部分由于这项工作的影响，Yoshua
Bengio获得了2018年图灵奖。

判别式方法似乎巧妙地绕过了分布学习中的一个根本困难。然而，严格来说，这种方法并没有完全解决这个根本困难。\cite{Goodfellow-2014}中表明，通过适当选择损失函数，极小极大公式在数学上等价于最小化\(\vX\)和\(\hat{\vX}\)之间的Jensen-Shannon距离（见\cite{Cover-Thomas}）。对于高维空间中的两个低维分布来说，这仍然是一个难题。因此，GAN通常依赖于许多启发式方法和工程技巧，并且经常遭受模式崩溃（mode
collapsing）等不稳定性问题的困扰。\footnote{尽管如此，这种极小极大公式提供了距离的实用近似。它简化了实现，并避免了直接计算距离时的一些注意事项。}

\paragraph{通过去噪和扩散进行生成。}
2015年，在GAN被引入并受到欢迎后不久，Surya
Ganguli和他的学生提出，由深度网络建模的迭代去噪过程可用于学习一般分布，例如自然图像的分布\cite{Sohl-Dickstein2015}。他们的方法受到了William
Feller在1949年研究的特殊高斯和二项式过程的性质的启发\cite{Feller1949OnTT}。\footnote{再次强调，这是在神奇的20世纪40年代！}

\begin{figure}[t]
    \centering
    \includegraphics[width=\linewidth]{\toplevelprefix/chapters/chapter1/figs/diffusion_pipeline.png}
    \caption{迭代去噪和压缩过程的图示，该过程从高斯分布\(q^{0} = \mathcal{N}(\vzero,
    \vI)\)开始，收敛到任意低维数据分布\(q^{L} = p\)。}
    \label{fig:diffusion}
\end{figure}

很快，基于分数函数（score
function）\cite{hyvarinen05a}的去噪算子（在\Cref{sec:denoising-intro}中简要介绍过）被证明更具一般性，并统一了去噪和扩散过程及算法\cite{song2019,song2020score,ho2020denoising}。\Cref{fig:diffusion}说明了通过执行一系列迭代去噪（或压缩）操作，将通用高斯分布\(q^{0}
= \mathcal{N}(\vzero, \vI)\)转换为（任意）经验分布\(p(\vx)\)的过程：
\begin{equation}
    \vz^{0} \sim  \mathcal{N}(\vzero, \vI) \xrightarrow{\hspace{2mm}
    g^{0} \hspace{2mm}} \vz^{1} \xrightarrow{\hspace{2mm} g^{1}
    \hspace{2mm}} \cdots \xrightarrow{\hspace{2mm} g^{L-1}
    \hspace{2mm}} \vz^{L} \sim p(\vx).
\end{equation}
到目前为止，去噪（和扩散）已经取代了GAN，成为学习图像和视频分布的主流方法，催生了Midjourney和Stability.ai等流行的商业图像生成引擎。
在\Cref{ch:denoising}中，我们将系统地介绍和研究用于学习一般低维分布的去噪和扩散方法。

\section{一种统一的方法}\label{sec:unifying-approach}
% Start of Selection
到目前为止，我们已经简要介绍了机器智能的主要目标和历史，以及与之相关的许多重要思想和方法。近年来，随着深度神经网络在经验上的成功，人们付出了巨大的努力来开发理论框架，以帮助我们理解经验设计的深度网络——无论是特定的看似必要的组件（例如，dropout
\cite{srivastava2014dropout}、归一化
\cite{ioffe2015batch,ba2016layer}、注意力机制
\cite{vaswani2017attention}），还是它们的整体行为（例如，双重下降
\cite{belkin2019reconciling}、神经坍缩 \cite{papyan2020prevalence}）。

部分受这一趋势的启发，本书追求几个重要且具有挑战性的目标：
\begin{itemize}
    \item 开发一个理论框架，对深度神经网络产生严格的数学解释。
    \item 确保学习到的数据分布的正确性以及与学习到的表征的一致性。
    \item 证明该框架能够带来高性能的架构，并指导进一步的实际改进。
\end{itemize}
近年来，越来越多的证据表明，通过利用前面讨论的经典分析低维模型的理论和解决方案（在\Cref{ch:classic-models}中有更深入的探讨），并整合相关领域（即信息/编码理论、控制/博弈论和优化）的基本思想，确实可以实现这些目标。本书旨在对这种新方法进行系统的介绍。

\subsection{学习简约表征}
\label{sec:computational-approach-compression}
任何学习任务成为可能的一个必要条件是，感兴趣的序列必须是\textit{可计算的}（computable），至少在艾伦·图灵（Alan
Turing）的意义上是如此\cite{Turing-1936}。也就是说，一个序列可以通过典型计算机上的程序来计算。\footnote{确实存在不可计算的明确定义的序列。}除了可计算之外，我们还要求计算是\textit{易处理的}（tractable）。\footnote{我们不需要考虑预测那些计算复杂性难以处理的事物——例如，随序列的长度或维度呈指数增长。}也就是说，学习和计算该序列的计算成本（空间和时间）不应随长度呈指数增长。此外，正如我们在自然界（以及机器智能的现代实践）中看到的那样，对于大多数实际任务，智能系统需要从非常高维空间（如视觉、声音和触觉）的海量数据中学习什么是可预测的。因此，对于智能而言，我们不需要考虑所有可计算和易处理的序列或结构；我们应该只关注那些允许其学习和计算算法具有\textit{可扩展}（scalable）实现的、可预测的序列和结构：
\begin{equation}
    \text{\textbf{可计算}} \;
    \Longrightarrow \; \text{\textbf{易处理}} \; \Longrightarrow \;
    \text{\textbf{可扩展}}.
\end{equation}

这是因为无论智能体使用什么算法来学习有用的信息，这些算法都必须是\textit{可扩展的}。更具体地说，算法的计算复杂性应该随数据的大小和维度优雅地扩展——通常是线性的甚至次线性的。在技术层面上，这要求算法赖以学习的操作只能利用可以从数据中高效计算出的预言机（oracle）信息。更具体地说，当维度很高且规模很大时，人们唯一能负担得起计算的预言机要么是关于数据的一阶几何信息\footnote{例如用线性子空间局部近似非线性结构并计算目标函数的梯度。}，要么是二阶统计信息\footnote{例如数据或其特征的协方差或相关性。正如我们将看到的，像Transformers这样流行的深度架构就依赖于此类信息。}。
本书的主要目标是开发一个理论和计算框架，在该框架内，我们可以系统地开发具有这种可扩展预言机和操作的高效且有效的解决方案或算法，以从采样数据中\textit{学习}低维结构，并随后学习预测函数。

\paragraph{通过压缩追求低维性。}
从我们在\Cref{sec:predictability}中给出的序列示例中可以清楚地看出，有些序列很容易建模和计算，而另一些则比较困难。序列的计算成本取决于预测函数\(f\)的复杂性。回归度\(d\)越高，计算成本就越高。函数\(f\)可以是一个简单的线性函数，也可以是一个任意难以指定和计算的非线性函数。

我们有理由相信，如果一个序列——无论我们选择什么度量标准——更难指定和计算，那么从其采样片段中学习它也会更加困难。然而，对于任何给定的可预测序列，实际上有许多（通常是无限多）种方法来指定它。例如，对于简单的序列\(x_{n+1}
= a x_{n}\)，我们也可以用\(x_{n+1} = a x_{n} + b x_{n-1} - b x_{n-1}\)来定义相同的序列。
因此，为任何给定的可计算序列开发一个客观且严格的“复杂性”概念将是非常有用的。

俄罗斯数学家安德雷·柯尔莫哥洛夫（Andrey
Kolmogorov）是最早为任何可计算序列定义复杂性的人之一。\footnote{许多人对序列复杂性的这一概念做出了贡献，最著名的是Ray
    Solomonoff和Greg
Chaitin。这三个人独立发展了算法信息论——Solomonoff在1960年，Kolmogorov在1965年\cite{Kolmogorov1998OnTO}，Chaitin在1966年左右\cite{Chaitin-1966}。}他提出，在计算同一序列的所有程序中，最短程序的长度衡量了其复杂性。这与奥卡姆剃刀原理相吻合——\textit{选择解释同一观察结果的最简单理论}。形式上，设\(p\)是在通用计算机\(\cU\)上生成序列\(S\)的程序。\(S\)的柯尔莫哥洛夫复杂性（Kolmogorov
complexity）为：
\begin{equation}
    ```latex
    K(S) = \min_{p\,:\, \cU(p) = S} L(p).
\end{equation}
因此，复杂度衡量了我们能够多么精简地指定或计算该序列。这个定义在概念上非常重要，并在历史上启发了计算复杂度和理论计算机科学领域的深刻研究。

最短的程序长度代表了序列的极限压缩，量化了我们从学习其生成机制中获得的收益。然而，柯尔莫哥洛夫复杂度通常是{\em 不可计算的}
\cite{Cover-Thomas}，并且难以近似。因此，它的实际用途很小——它既不能预测学习的难度，也不能评估我们学习的效果。

\paragraph{可计算的精简度量。}
出于实用目的，我们需要一种能够高效计算的复杂度度量，用于衡量由同一个预测函数生成的序列。\footnote{在实践中，我们通常关心的是学习预测函数
\(f\)，而不是由 \(f\) 生成的任何特定序列。} 请注意，柯尔莫哥洛夫复杂度不可计算的部分原因在于其定义是非构造性的。

为了引入一种可计算的复杂度度量，我们可以采取一种更具构造性的方法，正如克劳德·香农（Claude
Shannon）通过信息论框架所倡导的那样
\cite{Shannon-1948,Cover-Thomas}。\footnote{该框架已成功指导了通信行业80多年的工程实践。}
假设序列 \(S\) 提取自概率分布
\(p(S)\)，该分布的熵：\footnote{这里我们考虑的是微分熵，因为我们假设序列由连续变量组成。对于离散变量，我们将使用
\(H(S) = - \sum_{i}p(s_{i}) \log p(s_{i})\)。}
\begin{equation}
    h(S) \doteq -\int p(s) \log p(s) \odif{s}
    \label{eqn:entropy-definition}
\end{equation}
为其复杂度提供了一种自然的度量。这种度量也有一个自然的解释，即编码这样一个序列所需的平均二进制比特数，正如我们将在
\Cref{ch:denoising} 中看到的那样。

为了说明这种观点的主要思想，让我们取大量长序列片段：
\begin{equation}
    \{S_{1}, S_{2}, \ldots, S_{i}, \ldots, S_{N}\} \subset \R^{D},
\end{equation}
这些片段由预测函数 \(f\) 生成。不失一般性，我们假设所有序列具有相同的长度 \(D\)，因此每个序列可以被视为 \(\R^{D}\)
中的一个向量。我们引入一种编码方案（带有码本），记为 \(\cE\)，它将每个片段 \(S_{i}\) 映射到唯一的二进制比特流
\(\cE(S_{i})\)，以及相应的解码方案 \(\cD\)，它将二进制比特串映射回序列片段（假设误差为 \(\epsilon
\geq 0\)，因为这种编码方案可能是有损的）。最简单的方案是用 \(\epsilon\)-球填充由所有片段张成的空间，如
\Cref{fig:coding-schemes}
所示。然后我们对这些球进行编号，每个序列被编码为其最近的球的二进制索引，而每个二进制索引被解码回对应球的中心。因此，每个片段都可以从其比特流中以
\(\epsilon\) 的精度被恢复。
\begin{figure}
    \centering
    \includegraphics[width=0.5\linewidth]{\toplevelprefix/chapters/chapter1/figs/Coding-schemes.png}
    \caption{两种编码方案的比较。假设真实的数据分布位于两条带箭头的线上。来自这些线的样本可以用包含所有蓝色球的码本进行编码，或者仅用包含绿色球的码本进行编码。在相同精度下，第二种方案产生的编码长度/编码率要低得多。}
    \label{fig:coding-schemes}
\end{figure}

那么，对于编码 \(\cE\)，预测函数 \(f\) 的复杂度可以评估为它生成的所有序列 \(S\) 的平均编码长度
\(L(\cE(S))\)，这被称为编码率（coding rate）：\footnote{我们可以通过将 \(R(S \mid
\cE)\) 设为所有长度为 \(D\) 的片段 \(S\) 的期望编码长度，来使这一点更加精确。}
\begin{equation}
    R(S \mid \cE) = \Ex[L(\cE(S))] \approx \frac{1}{N}\sum_{i=1}^{N}
    L(\cE(S_{i})).
    \label{eqn:coding-rate}
\end{equation}
如果我们使用不同的编码方案（或码本），编码率度量将会改变。在实践中，我们对片段分布周围的低维结构了解得越清楚，我们就能设计出越高效的码本，如
\Cref{fig:coding-schemes} 所示。敏锐的读者可能会发现，从概念上讲，\Cref{fig:diffusion}
中说明的去噪过程非常类似于从使用蓝色球的编码方案过渡到使用绿色球的编码方案。

给定片段的两种编码方案 \(\cE_{1}\) 和 \(\cE_{2}\)，如果编码率的差值为正：
\begin{equation}
    R(S \mid \cE_{1}) -  R(S \mid \cE_{2}) > 0,
\end{equation}
我们可以说编码方案 \(\cE_{2}\) 更好。这种差异可以被视为一种度量，衡量了 \(\cE_{2}\) 比 \(\cE_{1}\)
多包含了多少关于数据分布的信息。在很大程度上，学习数据分布的目标等同于寻找最有效的编码和解码方案，从而在满足所需精度的前提下最小化编码率：
\begin{equation}
    \min_{\cE, \cD} R(S \mid \cE) \quad \text{subject to} \quad
    \operatorname{dist}(S, \cD(\cE(S))) \leq \epsilon.
\end{equation}
正如我们将在 \Cref{ch:lossy-compression} 中看到的，可达到的最小编码率与熵 \(H(S)\)
\eqref{eqn:entropy-definition} 的概念密切相关。

\begin{remark}\label{rem:computable-complexity}
    用显式编码方案衡量数据复杂度的视角，激发了几个旨在修正柯尔莫哥洛夫复杂度以获得更好可计算性的学习目标
    \cite{WallaceC1999}，包括1968年提出的最小消息长度（MML）\cite{WallaceC1968}
    和1978年提出的最小描述长度（MDL）\cite{Rissanen-1978,HansenM2001}。除了感兴趣的数据
    \(S\) 之外，这些目标通常还会计算编码方案 \(\cE\) 本身（包括其码本）的编码长度：\(L(\cE(S)) +
    L(\cE)\)。然而，如果目标是学习一个有限大小的码本并将其应用于大量序列，那么码本的摊销成本可以忽略不计，因为
    \begin{equation}
        \frac{1}{N}\left( L(\cE) + \sum_{i=1}^{N}
        L(\cE(S_{i}))\right) \approx \frac{1}{N}\sum_{i=1}^{N} L(\cE(S_{i}))
    \end{equation}
    当 \(N\) 变大时。
\end{remark}

同样，人们可以将由此产生的最优编码方案视为实现了对观测数据最佳压缩的方案。一般来说，与柯尔莫哥洛夫复杂度相比，任何编码方案给出的编码长度总是更大的：
\begin{equation}
    K(S) < L(\cE(S)).
\end{equation}
因此，最小化编码率/编码长度本质上是在最小化原本不可计算的柯尔莫哥洛夫复杂度的上限。

\subsection{学习信息丰富的表征}
如果目标仅仅是为了压缩给定数据本身，那么接近柯尔莫哥洛夫复杂度的最优代码在理论上将变得几乎是随机的或无结构的
\cite{Chaitin-1966}。\footnote{任何有结构的代码都可以被进一步压缩。} 然而，我们学习序列 $S$
的数据分布的真正目的是为了在未来的预测中轻松地重复使用它。因此，虽然压缩允许我们识别数据中的低维分布，但我们希望以一种\textit{结构化和有组织}的方式对这种分布进行编码，以便生成的表征信息丰富且使用高效\footnote{例如，在不同条件下对分布进行采样。}，作为一种“记忆”。\Cref{fig:expansion}
直观地说明了为什么这种转换是可取的。

\begin{figure}
    \centering
    \includegraphics[width=0.98\linewidth]{\toplevelprefix/chapters/chapter1/figs/coding-transform.png}
    \caption{将识别出的低维数据分布转换为更具信息量和结构化的表征。}
    \label{fig:expansion}
\end{figure}
正如我们将在 \Cref{ch:lossy-compression}
中展示的那样，通过选择一种基于所选编码方案编码率的\textit{信息增益}的自然度量，可以精确地促进最终表征中这些期望的结构。\footnote{信息增益是算法信息论中一个普遍存在的概念，并已被纳入该主题的几种方法中，最近的是
    \(\mathcal{V}\)-信息
\cite{xu2020theory}。与我们讨论过的许多其他复杂度度量类似，这种特定的实例化随着数据支撑集变得更加紧凑和低维，会变得越来越病态，因此不适合我们的需求。}
在本书中，这种显式且具构造性的编码方法为学习现实世界数据中低维结构的良好表征提供了一个强大的计算框架。在许多具有实际重要性的情况下，编码长度函数可以被高效计算或精确近似。在一些条件良好的情况下，我们甚至可以获得闭式公式——例如，子空间和高斯模型。

此外，这种计算框架引出了一种原则性的方法，自然地揭示了深度网络在这一学习过程中所扮演的角色。正如我们将在
\Cref{ch:deep-networks}
中系统推导的那样，深度网络的各层执行的操作是在增量式地优化感兴趣的目标函数。从这个角度来看，深度网络的作用可以被精确地解释为模拟某种迭代优化算法——例如梯度下降——以优化信息增益的目标。由此产生的深度架构的各层可以被赋予精确的统计和几何解释，即执行增量式的压缩编码和解码操作。因此，推导出的深度网络变成了在数学上完全可解释的透明“白盒”。

\subsection{学习一致的表征}

\label{sec:consistency}

为了总结我们到目前为止的讨论，让我们将数据记为
\begin{equation}
    \vX = \{S_{1}, S_{2}, \ldots, S_{i}, \ldots, S_{N}\} \subset \R^{D},
\end{equation}
并令 \(\vZ = \cE(\vX)\) 为 \(\vX\) 通过某个编码器 \(\cE\) 得到的代码：
\begin{equation}
    \vX  \xrightarrow{\hspace{2mm} \cE\hspace{2mm}} \vZ.
    \label{eqn:open-loop}
\end{equation}
在机器学习的语境中，\(\vZ\) 通常被称为 \(\vX\) 的“特征”。请注意，在不知道 \(\vX\)
潜在分布的情况下，我们不知道应该使用哪个编码器 \(\cE\) 来保留关于 \(\vX\)
分布的最有用信息。在实践中，人们通常从尝试一种能很好地服务于特定任务的紧凑编码方案开始。特别是，他们试图学习一个编码器，该编码器优化了所学表征的某种（经验）精简度量：
\begin{equation}
    \min \rho(\vZ).
\end{equation}

\begin{example}\label{example:image_classification_feature_learning}
    例如，图像分类就是这样一种情况：我们将同一类别的所有图像分配给一个单一的代码，将不同类别的图像分配给不同的代码，例如“独热（one-hot）”向量：
    \begin{equation}
        \vx \mapsto \vz \in \{  [1, 0, 0, \ldots , 0, 0], \;  [0, 1,
        0 \ldots, 0, 0], \; \ldots, \;  [0,0,0, \ldots, 0, 1].\}
        \label{eqn:class-labels}
    \end{equation}
    现在，分类器 \(f(\cdot)\) 可以被建模为一个函数，它预测给定 \(\vx\) 属于 \(k\)
    个类别中每一个的概率：\(\hat{\vz} = f(\vx) \in
    \R^{K}\)。那么分类器的“好坏”可以通过所谓的{\em 交叉熵}来衡量：\footnote{交叉熵可以被视为 \(\vz\)
        的真实分布与预测 \(\hat{\vz}\) 的分布之间的距离度量。如果我们使用 \(\hat{\vz}\) 的最优码本来编码
        \(\vz\)，它也可以被视为 \(\vz\) 的期望编码长度。当 \(\vz\) 和 \(\hat{\vz}\)
    具有相同的分布时，交叉熵达到其最小值。}
    \begin{equation}
        \ell(\hat{\vz}, \vz) = \sum_{k=1}^{K} - z_{k} \log \hat{z}_{k},
    \end{equation}
    其中 \(z_{k}\) 表示向量 \(\vz\) 的第 \(k\) 个分量。正如深度网络的早期实践所表明的那样
    \cite{krizhevsky2012imagenet}，如果给定足够的数据，这种编码方案通常可以由深度网络表示，并通过优化交叉熵以端到端的方式进行学习。
\end{example}

交叉熵损失 \(\ell(\hat{\vz}, \vz)\) 可以被视为与适合分类的特定编码方案族相关联的一种特殊的精简度量
\(\rho(\vz)\)。然而，这种编码显然是\textit{非常有损的}。除了类别类型之外，学习到的 \(\vz\) 不包含关于
\(\vx\) 的任何其他信息。例如，通过为图像分配（代表）类别标签“苹果”的代码，我们仅从标签中不再知道原始图像中是哪种特定类型的苹果。

当然，另一个极端是要求编码方案是\textit{无损的}。也就是说，在 \(\vx\) 及其代码 \(\vz\)
之间存在一一映射。然而，正如我们将在 \Cref{ch:lossy-compression} 中看到的，除非 \(\vx\)
是离散的，否则无损编码（或压缩）是不切实际的。对于连续随机变量，我们只能考虑有损编码方案，以便数据的编码长度可以是有限的。也就是说，我们只将数据编码到某个规定的精度。正如我们将在
\Cref{ch:lossy-compression}
中进一步阐述的那样，有损编码不仅仅是一种实际选择；它在使从分布的有限样本中学习潜在连续分布成为可能方面发挥着基础性作用。

对于许多学习目的，我们希望特征 \(\vz\) 尽管是\textit{有损的}，但能保留比仅仅是类别类型更多的关于 \(\vx\)
的信息。在本书中，我们将引入一种更通用的精简度量，该度量基于与更通用的编码方案族（即使用子空间混合或高斯混合进行编码）相关联的编码长度/编码率。该族有能力在一定精度内紧密逼近任意现实世界的分布。正如我们将在
\Cref{ch:lossy-compression} 和 \Cref{ch:deep-networks}
中看到的，这种度量不仅将保留关于 \(\vX\) 分布的大部分信息，而且还将促进学习到的表征 \(\vZ\) 形成某些期望的几何和统计结构。

\paragraph{用于一致性的双向自编码。}
在更广泛的学习背景下，压缩编码方案 \(\cE\) 的主要目标是识别数据 \(\vX\)
中的低维结构，以便它们可用于预测原始数据空间中的事物。这要求学习到的编码方案 \(\cE\) 允许一个高效的解码方案，记为
\(\cD\)，它将特征或编码数据 \(\vZ\) 映射回数据空间：
\begin{equation}
    \vX   \xrightarrow{\hspace{2mm} \cE\hspace{2mm}} \vZ
    \xrightarrow{\hspace{2mm} \cD \hspace{2mm}} \hat{\vX}.
    \label{eqn:auto-encoding}
\end{equation}
我们将这样的编码-解码对 \((\cE, \cD)\)
称为\textit{自编码（auto-encoding）}。\Cref{fig:autoencoder} 说明了这一过程。
\begin{figure}
    \centering
    \includegraphics[width=0.7\linewidth]{\toplevelprefix/chapters/chapter1/figs/Autoencoder.jpg}
    \caption{自编码器架构示意图。}
    \label{fig:autoencoder}
\end{figure}

理想情况下，解码近似于编码的“逆”，因此从 \(\vZ\) 解码出的数据（分布）\(\hat{\vX}\)
在某种程度上与原始数据（分布）\(\vX\) 相似。\footnote{我们稍后将明确“相似”的含义。} 如果这一点成立，我们就可以从
\(\vZ\) 中恢复或预测原始数据空间中发生的事情。在这种情况下，我们说该对 \((\cE, \cD)\)
提供了一个\textit{一致的（consistent）}自编码。对于大多数实际目的，我们不仅需要这样的解码存在，还需要它能够被高效地实现并在物理上可执行。例如，如果
\(\vx\) 是实值的，那么任何解码方案要在有限状态机上实现都需要量化，正如我们将在
\Cref{ch:lossy-compression}
中解释的那样。因此，一般来说，可实现的编码和解码方案必然是有损的。一个核心问题是如何学习一个紧凑的（有损）表征
\(\vZ\)，它仍然能够很好地预测 \(\vX\)。

一般来说，编码器和解码器都可以由深度网络建模和实现，并通过求解以下形式的优化问题来学习：
\begin{equation}
    \min_{f, g} [\ell(\vX, \hat{\vX}) + \rho(\vZ)], \qquad
    \text{where} \quad \vZ = f(\vX), \quad \hat{\vX} = g(\vZ)
    \label{eqn:auto-encoding-objective}
\end{equation}
其中 \(\ell(\cdot, \cdot)\) 是一个促进 \(\vX\) 和 \(\hat{\vX}\)
之间相似性的距离函数\footnote{取决于 \(\ell\) 的选择，可以是样本级别的或分布级别的。}，而 \(\rho(\vZ)\)
是一个促进 \(\vZ\) 的精简性和信息丰富度的度量。经典的主成分分析（PCA）\cite{JolliffeI2002}
是一致自编码的典型例子，我们将在 \Cref{ch:classic-models} 中对其进行详细研究，作为更一般低维结构的先导。在
\Cref{ch:consistent-representations} 中，我们将研究如何为一般的（例如非线性的）低维分布学习一致的自编码。

\subsection{学习自洽的表征}
在上述自编码目标中，必须评估解码数据 \(\hat{\vX}\) 与原始数据 \(\vX\)
的接近程度。这通常需要外部监督或对适当相似性度量的了解。计算 \(\hat{\vX}\) 和 \(\vX\)
之间的相似性可能极其昂贵，甚至是不可能的。\footnote{例如，最小化两个随机变量之间的分布距离在统计和计算上都是困难的。}
在自然界中，动物自主学习，而不会在数据空间中将它们的估计 \(\hat{\vX}\) 与真实值 \(\vX\)
进行比较；它们通常缺乏这种选择。\footnote{在动物和人类中，眼睛在视觉数据到达大脑之前很久就对其进行了处理；大脑没有处理原始视觉数据的机制，必须依赖眼睛的特征来进行任何学习。}

那么，一个系统如何在没有外部监督或比较的情况下进行学习呢？它如何在没有直接比较的情况下知道 \(\hat{\vX}\) 与 \(\vX\)
是一致的呢？这引出了“闭环（closing the loop）”的想法。在温和的条件下（在
\Cref{ch:consistent-representations} 中有精确描述），确保 \(\vX\) 和
\(\hat{\vX}\) 之间的一致性可以简化为将 \(\hat{\vX}\) 编码为 \(\hat{\vZ}\)，并检查
\(\vZ\) 和 \(\hat{\vZ}\) 之间的一致性。我们将这个概念称为\textit{自洽性（self-consistency）}，如下所示：
\begin{equation}
    \vX   \xrightarrow{\hspace{2mm} \cE\hspace{2mm}} \vZ
    \xrightarrow{\hspace{2mm} \cD \hspace{2mm}} \hat{\vX}
    \xrightarrow{\hspace{2mm} \cE\hspace{2mm}} \hat{\vZ},
    \label{eqn:closed-loop}
\end{equation}
我们将这个过程称为\textit{闭环转录（closed-loop
transcription）}\footnote{受DNA与RNA或其他蛋白质之间转录的启发。}，如 \Cref{fig:closed-loop} 所示。

\begin{figure}[t]
    \centering
    \includegraphics[width=0.9\linewidth]{\toplevelprefix/chapters/chapter1/figs/diagrams_redu_gan_2.pdf}
    \caption{闭环转录示意图。这里我们使用映射 \(f\) 来表示编码器 \(\cE\)，使用 \(g\) 来表示解码器
    \(\cD\)。}  \label{fig:closed-loop}
\end{figure}

可以说，任何自主的智能体只需要学习观测数据 \(\vX\) 的自洽表征
\(\vZ\)，因为在原始数据空间（通常意味着在外部世界）中检查一致性要么太昂贵，要么在物理上甚至不可行。闭环公式允许人们通过一个仅依赖于内部（学习到的）特征
\(\vZ\) 的最小-最大博弈（min-max game）来学习最优的编码 \(f(\cdot, \vtheta)\) 和解码
\(g(\cdot, \veta)\)：
\begin{equation}
    \max_{\vtheta}\min_{\veta} [\ell( \vZ, \hat{\vZ}) + \rho(\vZ)],
    \qquad \text{where} \quad \vZ = f(\vX), \quad \hat{\vX} = g(\vZ),
    \quad \hat{\vZ} = f(\hat{\vX})
    \label{eqn:closed-loop-objective}
\end{equation}
其中 \(\ell( \vZ, \hat{\vZ})\) 是一个基于特征 \(\vZ\) 和 \(\hat{\vZ}\)
编码率的损失函数，正如我们将看到的，它可能更容易计算。这里同样，\(\rho(\vZ)\) 是某种促进 \(\vZ\)
的精简性和信息丰富度的度量。有些令人惊讶的是，正如我们将在 \Cref{ch:consistent-representations}
中看到的，在相当温和的条件下（例如 \(\vX\) 足够低维），\((\vZ, \hat{\vZ})\) 之间的自洽性意味着
\((\vX, \hat{\vX})\)
之间的一致性！此外，我们还将看到，闭环系统将允许我们以\textit{连续和增量}的方式学习分布\footnote{也就是说，一次学习一个类别，甚至一次学习一个样本。}，而不会遭受与开环模型相关的灾难性遗忘等问题。

% Start of Selection
\section{连接机器智能的理论与实践}
到目前为止，我们已经介绍了几个相关的框架，用于为给定的数据分布 \(\vX\) 学习紧凑且结构化的表征 \(\vZ\)：
\begin{itemize}
    \item 开放式编码 \eqref{eqn:open-loop}；
    \item 双向自编码 \eqref{eqn:auto-encoding}；
    \item 闭环转录 \eqref{eqn:closed-loop}。
\end{itemize}
在本书中，我们将依次系统地研究这三个框架：
\begin{equation}
    \text{\textbf{开放式}} \; \Longrightarrow \;
    \text{\textbf{双向}} \; \Longrightarrow \;
    \text{\textbf{闭环}},
\end{equation}
分别在 \Cref{ch:deep-networks}、\Cref{sec:consistent-representation} 和
\Cref{sec:self-consistency} 中进行探讨。

正如我们将在本书中看到的，这些框架（包括本章前面讨论的许多框架）现在可以被视为试图{\em 部分地}解决学习数据 $\X$ 分布的{\em
完整且可计算}表征的问题：
\begin{equation}
    \X \longleftrightarrow \Z \longleftrightarrow \boldsymbol{G},
\end{equation}
这包括两组映射：
\begin{enumerate}
    \item 数据 $\X$ 及其信息丰富的特征表征 $\Z$ 之间的自编码，
    \item 学习到的表征 $\Z$ 与简单分布 $\boldsymbol{G}$（例如高斯分布）之间的扩散和去噪过程。
\end{enumerate}
后一个过程允许我们访问和采样学习到的结构化低维特征 $\Z$，用于重新生成原始数据 $\hat \X$
等任务。\footnote{请注意，即使学习到的特征位于单个低维线性子空间上，识别该子空间最有效的方法本质上也是通过去噪过程：我们将在第
\ref{ch:classic-models} 章中详细阐述的幂迭代（power iteration）。}
与这两个（转换）过程相关的整体几何图像如图 \ref{fig:mental-picture} 所示。

正如我们将在第 \ref{ch:deep-networks}
章中看到的，现代深度网络的作用可以被精确地解释为实现上述分布之间转换所需的算子\footnote{正如我们将看到的，通过展开优化（unrolled
optimization）}。凭经验开发出的流行深度神经网络架构（包括 CNN、ResNet 或
Transformer）将通过这种方法在数学上得到解释、证明和改进。

\begin{figure}[t]
    \centering
    \includegraphics[width=0.6\linewidth]{\toplevelprefix/chapters/chapter1/figs/book-cover-2.jpg}
    \caption{数据 $\X$ 学习到的信息丰富且结构化的特征表征 $\Z$ 示意图，该表征可以通过从
    $\boldsymbol{G}$ 开始的去噪过程被高效访问。}  \label{fig:mental-picture}
\end{figure}

近年来，人们提出了许多理论框架来帮助理解深度网络。然而，大多数框架未能提供可扩展的解决方案，以在现实世界的数据和任务上匹配经验方法的性能。许多理论对进一步的经验改进几乎没有提供实际指导。\Cref{ch:conditional-inference}
和 \Cref{ch:applications}
将展示本书提出的框架如何弥合理论与实践之间的差距。\Cref{ch:conditional-inference}
将展示如何使用学习到的分布及其表征，为涉及（条件）生成、估计和预测的实际任务进行（贝叶斯）推断。\Cref{ch:applications}
将提供令人信服的实验证据，证明从第一性原理设计的网络和系统可以在视觉表征学习、图像分类、图像补全、分割和文本生成等任务上实现具有竞争力——甚至可能更优越——的性能。

\paragraph{回到智能。}
正如开头所提到的，任何智能体的一项基本任务是从感知数据中学习可预测的信息。我们现在理解了这项任务的计算本质，并认识到这是一个永无止境的过程，原因有二：
\begin{itemize}
    \item 从数据中学习到的记忆或知识（例如通过编码和解码方案）不太可能是正确或最优的。当对新观测的预测包含错误时，智能必须能够自我改进。
    \item 观测数据尚未涵盖所有可预测的信息。智能必须能够识别当前的记忆或知识何时不足，并在新信息出现时随时获取。
\end{itemize}

因此，智能\textit{并非}指提前收集所有信息数据并训练模型来记忆其（低维）分布。相反，它需要能够不断改进当前对分布的知识并在需要时获取新信息的计算机制。任何智能体或系统——无论是动物、人类、文明，还是整个科学界——的一个基本特征是，\textit{能够通过关于其当前记忆或知识在解释或预测观测结果方面表现如何的闭环反馈，自主地不断改进或获取新信息或知识}。因此，在某种程度上，我们可以将智能与知识\footnote{这里的“知识”是最广义上的，既包括经验记忆，也包括科学知识。}之间的关系符号化地表示为：
\begin{eqnarray*}
    \operatorname{Intelligence}(t) &=&
    \odv*{\operatorname{Knowledge}}{t}(t), \\
    \operatorname{Knowledge}(t)  &=& \int_{0}^{t}
    \operatorname{Intelligence}(s) \odif{s}.
\end{eqnarray*}
闭环反馈框架是任何具备自我改进和自我学习能力的系统的通用机制——强化（如自然选择中的强化）或博弈是其原始形式，而科学探究通过假设或猜想以及证伪或验证代表了其最高级的形式。自然界中所有智能体和系统在各个层次和尺度上都采用这种闭环机制进行学习。这种普遍性启发了早期用自主机器或半自主计算机对智能进行建模的尝试，特别是诺伯特·维纳（Norbert
Wiener）在20世纪40年代开创并倡导的控制论（Cybernetics）运动。

我们希望本书能帮助读者更好地理解智能背后的科学目标、数学原理和计算机制，至少是形成有用记忆或经验知识的那部分。我们相信，这项工作不仅架起了理论与实践的桥梁，也连接了现在与未来：它阐明了我们迄今为止所做的工作，同时也揭示了我们尚未完成的工作。因此，它为未来研究人类独有的更高层次智能背后可能的新机制奠定了基础——即1956年达特茅斯会议所指的真正的“人工智能”。我们认为，这种机制能够创造抽象概念并进行演绎推理，超越了本书所研究和解释的、以及当前“AI”技术正在利用的基于经验数据压缩记忆和进行归纳推理的范畴。我们将在最后的\Cref{ch:future}中讨论关于理解并潜在实现更高级智能的开放性问题。

\end{document}
